% =====================================================================
%  Harald Høffding, ERINDRINGER. København: Gyldendalske Boghandel,
%  Nordisk Forlag, 1928 (MCMXXVIII). First printing ("Oplag: 2000
%  Eksemplarer"; Gyldendals Forlagstrykkeri). 324 pp. + Indhold.
%
%  THIS FILE IS A SKELETON WITH SELECTED SECTIONS TRANSCRIBED.
%  Most of the book is NOT yet transcribed: every chapter stands here as
%  its printed head and a "% [text to be added: pp. X--Y]" marker, and
%  each gap is announced in the PDF by a visible \gap line.  What has
%  been transcribed, and how its accuracy was established, is stated in
%  section 5 below.  Nothing in this file outside those ranges has been
%  read against the page.
%
%  Diplomatic transcription from the scan
%  `bibliotek/Høffding, Harald/1928-erindringer.pdf`, 327 PDF pages,
%  376.55 x 502.067 pt, each text page one 600 ppi greyscale image with an
%  ABBYY FineReader 14 text layer (download cover sheet from Danskernes
%  Historie Online / Slægtsforskernes Bibliotek), sha256
%  ae8d8ab9445ea03d8d50a7aaab96a4d0090f5e5fdd2945d4f2e88413fddbe059.
%
%  WITHDRAWN PREDECESSOR.  Until 23 Sep 2026 this directory held a text
%  of selected paragraphs of chapter II derived from the Lindhardt og
%  Ringhof e-book (2021) of the 1928 text, not from the page, carrying 36
%  page markers invented for a source that has no pagination, together
%  with an English translation of that text.  Both were withdrawn and
%  are kept on disk only, untracked, as
%  transcription.tex.bak.ebook-withdrawn-20260923 and
%  translation.tex.bak.ebook-withdrawn-20260923.  THE PRESENT TEXT IS NOT
%  DERIVED FROM THEM: it was built from scratch at the scan, and neither
%  they nor the e-book were opened by anyone transcribing or collating
%  it.  Its page markers were read at the image.
%
%  This header is the edition's apparatus.  The harness (check.py,
%  splice.py, joints.py, the batch brief, RESUME-NOTES.md) is not tracked
%  in git and is unrunnable without the scan, so everything an outside
%  reader needs in order to check this text against the scan is here.
%
% ---------------------------------------------------------------------
%  1. PAGE MAP  (established 23 Sep 2026)
% ---------------------------------------------------------------------
%
%      printed p. N  =  PDF page N + 2, uniform       (pp. 7--324 = PDF 9--326)
%
%  Measured from the scan's own folios (paginate.measure_offset: offset 2
%  from 216 folios, unanimous, no second value anywhere in the book, so
%  no plate or unnumbered leaf interrupts the sequence after PDF 9), and
%  read at the image on PDF 12, 61, 97, 150, 240, 315 (folios 10, 59, 95,
%  148, 238, 313), and again by every batch transcriber on every page
%  transcribed.  Folios are centred at the head in old-style figures
%  (which the text layer reads as letters: "IO" for 10, "I2J" for 121).
%  Chapter-opening pages carry no folio.  There is no running head.
%
%  Leaves (PDF pages):
%      PDF   1     download cover sheet (not part of the book)
%      PDF   2     library binding, typed label "Harald Høffding / Erindringer"
%      PDF   3     original wrapper (silhouette; GYLDENDALSKE BOGHANDEL)
%      PDF   4     half-title, ERINDRINGER
%      PDF   5     portrait photograph with facsimile signature
%      PDF   6     TITLE PAGE
%      PDF   7     title verso: COPYRIGHT 1928 BY GYLDENDALSKE / BOGHANDEL
%                  NORDISK FORLAG - COPENHAGEN / [short rule] / OPLAG:
%                  2000 EKSEMPLARER / [short rule] / PRINTED IN DENMARK /
%                  GYLDENDALS FORLAGSTRYKKERI / KØBENHAVN
%      PDF   8     FORORD, unnumbered (p. [6] by sequence), dated
%                  "Carlsberg, 14. Marts 1928."             [not transcribed]
%      PDF   9--326  text, pp. 7--324
%      PDF 327     INDHOLD, unnumbered (p. [325] by sequence), with a
%                  five-line note beneath the list on when the sections
%                  were written                                [transcribed]
%
% ---------------------------------------------------------------------
%  2. STRUCTURE  (from the images, not the OCR)
% ---------------------------------------------------------------------
%
%  Eight chapters.  Each opens on a fresh page with a centred head that is
%  the Roman numeral ALONE -- no full stop, no title, no dates, no rule --
%  and is set \chapter*{I} etc., as printed.  The year-ranges exist only
%  in the Indhold; the table-of-contents entries and the "pp." ranges in
%  the \gap lines are EDITORIAL.  Only chapter I opens with a dropped
%  initial (a two-line N).
%
%      I     1843--1861   pp.   7--37        V     1894--1904   pp. 175--223
%      II    1861--1874   pp.  38--93        VI    1904--1914   pp. 224--265
%      III   1874--1883   pp.  94--120       VII   1914--1922   pp. 266--310
%      IV    1883--1894   pp. 121--174       VIII  1922--1924   pp. 311--324
%
%  Chapter IV: the Indhold prints "121" (read at 600 dpi: old-style
%  1-2-1; the final figure carries a small nick at its foot, which is
%  what the text layer read as "J").  Chapter IV's head does stand at the
%  top of p. 121 (PDF 123), and p. 120 ends chapter III.  So III is pp.
%  94--120 and IV pp. 121--174.
%
%  Chapter VI's head (p. 224, PDF 226) PRINTS "V": the I did not ink,
%  leaving a faint stub beside the V (read at the image twice).  It is
%  transcribed as printed, with a comment.  The Indhold prints "VI·" with
%  a raised point.
%
%  Chapter ends, as seen at reconnaissance (NOT yet transcribed or
%  collated): p. 265 "Sluttet d. 21. December 1916." (body type); p. 310
%  "Sluttet d. 15. Marts 1922." (small, set right); p. 324 "Den 10. Juli
%  1924." (small) and a note "¹) Vilhelm Thomsen døde 12. Maj 1927.
%  [Senere Anmærkning]."  These datelines, and the note printed beneath
%  the Indhold (transcribed and collated), bear on how the book is read:
%  by the Indhold's note, chapters I--VI were written in the winter of
%  1915--16, VII in 1922 and VIII in the spring of 1924, and a
%  read-through before printing added parentheses and notes, which are
%  part of the 1928 text and are transcribed as such.  (Chapter VI's own
%  dateline, December 1916, is the book's reading, recorded not
%  reconciled.)
%
% ---------------------------------------------------------------------
%  3. TYPOGRAPHY AND TRANSCRIPTION CONVENTIONS
% ---------------------------------------------------------------------
%
%  Body in ANTIQUA (roman) with old-style figures, not Fraktur.
%  Orthography as printed: aa (å only in Swedish words), capitalised
%  nouns, 1928 spelling.  Printer's errors are transcribed AS PRINTED and
%  logged in a % comment at the site.
%
%  Emphasis in this book is ITALIC, not letterspacing: a width survey of
%  every word found no letterspaced word in the body (calibration at 300
%  dpi, p. 66: inter-letter gaps 2--4 px in italic and roman alike,
%  21--25 px between words).  Italic -> \textit{}.  Letterspacing, where
%  it occurs (capital heads such as INDHOLD), -> \emph{}.
%
%  Quotation marks „...“ as printed; apostrophe ’ ("Hegel’s").  Ranges
%  and dashes as printed with an em rule (1865---66).  Footnotes are
%  marked "¹)" and numbered afresh on each printed page; they are set
%  \footnote[N]{...} with the printed number, so the numbering follows
%  the page and not LaTeX.  The quoted letters of pp. 63--69 are set in
%  the book in body type as ordinary paragraphs, without quotation marks,
%  with a white line before and after the run; so here.  A white line in
%  the print -> \whiteline; a short centred rule -> \secrule.  Spaced
%  ellipses ("Fantasien ... Det") are as printed.
%
%  \supplied{...} (printed ⟨...⟩) marks letters SUPPLIED BY THE EDITOR
%  where the SCAN, not the print, has lost them.  One instance: p. 61,
%  "⟨tæn⟩kte" (see the comment there and section 5).
%
%  Line-break hyphens are resolved.  A page turn inside a word is marked
%  inside the word: the first half ends its line with "%", the page
%  comment follows, and \opage{N} stands before the second half, so the
%  word is joined and the marginal number falls at the turn (e.g.
%  "Generalforsamlin" / \opage{42}"ger", pp. 41/42).  Where a transcribed
%  range begins or ends inside a word whose other half lies on an
%  untranscribed page, the half on the page is given as printed (with its
%  hyphen at an end) and a comment says so.  Folios, signature
%  marks at the page foot ("Erindringer" + sheet number, "5*"), and
%  library matter are not transcribed.
%
%  Every printed page transcribed carries, on its own line, a comment
%  "% --- p. N ---", and \opage{N} at the exact word the page begins
%  with.  All \opage markers were read at the image; there are no \apage
%  markers in this file.
%
% ---------------------------------------------------------------------
%  4. WHAT IS TRANSCRIBED  (and what is not)   [23 Sep 2026]
% ---------------------------------------------------------------------
%
%  Transcribed from the image, and collated (section 5):
%
%    pp. 39--51  (ch. II)  the student years and the teachers.  Opens
%                mid-sentence at the head of p. 39; ends mid-sentence
%                at the foot of p. 51 ("Forudsætning,").
%    pp. 55--69  (ch. II)  opens mid-word at the head of p. 55
%                ("kommende", the word begun on p. 54).  The press
%                discussion of winter 1865--66 over Rasmus Nielsen's
%                doctrine of faith and knowledge opens at the head of
%                p. 59.  The letters, in chronological order, begin on
%                p. 63 and run, with Høffding's parenthesis of pp. 67--68,
%                to p. 69 ("26. Juni 1870"); the range was carried to
%                p. 69 so that the whole run of letters stands.  Ends
%                mid-word at the foot of p. 69 ("efterhaan-").
%    pp. 77--82  (ch. II)  opens mid-word ("tede", begun on p. 76); ends
%                mid-sentence ("De fandt i nogle Aar").
%    pp. 90--95  (ch. II/III)  opens mid-sentence; chapter II ends on
%                p. 93; chapter III's head and first pages, pp. 94--95;
%                ends mid-sentence ("om").
%    Indhold, p. [325] (PDF 327), with the note beneath it.
%
%  The title page (PDF 6) is set from a reading at the image, checked
%  once more at the image for wording, separators (hyphens) and line
%  division; it was not collated word by word.
%
%  NOT transcribed -- skeleton only, each gap marked by a \gap line and
%  a "% [text to be added: pp. X--Y]" marker: the Forord (p. [6]);
%  chapter I (pp. 7--37); in chapter II pp. 38, 52--54, 70--76, 83--89;
%  chapter III pp. 96--120; chapters IV--VIII (pp. 121--324).  That is
%  some 280 of the book's 318 text pages.
%
% ---------------------------------------------------------------------
%  5. ACCURACY  (transcribed ranges only)   [23 Sep 2026]
% ---------------------------------------------------------------------
%
%  Everything in this section concerns the 41 pages listed in section 4
%  (40 text pages and the Indhold).  NOTHING ELSE IN THIS FILE HAS BEEN
%  TRANSCRIBED, so nothing else has an accuracy to state.
%
%  Word-accuracy: every batch diffed against the scan's embedded ABBYY
%  layer before splicing (ocrdiff.py --frag, offset 2).  7 batches
%  (pp. 39--45, 46--51, 55--61, 62--69, 77--82, 90--95, Indhold),
%  4 candidates, 0 transcription errors corrected, 4 the witness's fault,
%  0 unresolved.  See TRANSCRIPTION-PLAYBOOK.md §3 step 4.  The tool was
%  calibrated on these fragments: one short word (og, at, i, som, ...)
%  dropped at random mid-page on each of 40 pages was flagged on 34.
%  And because the transcribers took their words from that layer, the
%  diff cannot see a layer error they copied.  The readings below exist
%  for those two classes.
%
%  Method.  (1) TRANSCRIPTION: one agent per batch, words from the text
%  layer corrected at the image (300 dpi; 600 dpi, the scan's native
%  resolution, for anything in doubt); structure, italics, punctuation,
%  every figure and every footnote from the image; folio read on every
%  page.  (2) SECOND READING of ALL 41 pages, word by word and mark by
%  mark at the image, by four readers who had transcribed nothing and
%  were forbidden the text layer: 0 discrepancies in words, letters,
%  figures, punctuation, italics, paragraphing, footnotes or page
%  markers; every printer's error logged by the transcribers confirmed
%  at the image; three marks left undecidable (pp. 42, 44, 66).
%  (3) THIRD READING by a fifth reader who had done neither: (a) a BLIND
%  CONTROL -- pp. 57, 60, 64, 65 (seeded draw, 20260923) transcribed from
%  the image before the reader saw this file, then compared token by
%  token (words and single marks): 1,626 tokens, 0 differences, italics
%  identical; (b) the disputes settled at 600 dpi: p. 66 the comma after
%  "Kulturformer" is italic (CORRECTED from roman); pp. 42 (end of
%  note 3) and 44 ("Testamente.") points, probable (kept); p. 61 see
%  below; pp. 41, 94 and the Indhold's "VI·" and "121" confirmed.
%  The e-book text, the withdrawn files and every other edition were
%  closed to all five readers and all seven transcribers.
%
%  Result: 0 word errors found in 41 pages; 1 mark corrected (p. 66,
%  italic comma).  Stated bounds, not boasts: collation that found
%  nothing in 41 pages puts the word-error rate below about 0.07 per
%  page (95%, rule of three) ONLY IF collation catches every error; the
%  blind control, which does not depend on that, found 0 differences in
%  1,626 tokens (95% upper bound about 0.2% of tokens).
%
%  SUPPLIED TEXT.  p. 61: the left ends of the page's last two lines are
%  blank in this scan (the clean edge suggests the corner was covered in
%  scanning; paper loss cannot be excluded from this scan alone).
%  "Men" (second line) is read from a surviving part of the M and "en".
%  In "⟨tæn⟩kte" (first line) only the top of the k, a trace of the n and
%  "kte" survive; "tæn" is the editor's conjecture from the lost width
%  and the sense, set \supplied{tæn}.  TO CHECK against another copy.
%
%  Printer's errors transcribed as printed and logged at the site:
%  p. 41 "andelig"; p. 41 "Kierkegaard „Til Selvprøvelse“" (no genitive
%  s); p. 44 "sprængtes.Hinsidighedens"; p. 46 "mine legemlige
%  Konstitution"; p. 59 "Rasmus Nielsens hørte"; p. 60 "dennne";
%  p. 60 "( som"; p. 63 "jeg ngen Præst"; p. 78 "sine unge ... Hustru";
%  p. 78 "Tillidt"; p. [325] "VI·".  Noted as printed without being
%  called errors: p. 47 "ligened"; p. 55/56 "Mariboes"/"Maribos";
%  p. 58 "spasere"; p. 59 "dér"; p. 61 a risen space between "og" and
%  "Brøchners" (not typed); p. 63 a letter's date line without its em
%  rule;
%  p. 94 "drejede sig om" without "det"; p. [325] "Paranteser".
%  Quote balance: „ 64 / “ 64.
%
% =====================================================================

\documentclass[11pt]{book}
\usepackage[utf8]{inputenc}
\usepackage[T1]{fontenc}
\usepackage{libertinus}
\usepackage{libertinust1math}
\usepackage{textalpha}
\usepackage{amsmath}
\usepackage[danish]{babel}
\usepackage{setspace}
\usepackage{fancyhdr}
\setlength{\headheight}{28pt}
\addtolength{\topmargin}{-13.5pt}
\pagestyle{fancy}
\fancyhf{}
\fancyhead[LE,RO]{\thepage}
\fancyhead[RE]{\textit{Harald Høffding: Erindringer}}
\fancyhead[LO]{\textit{\leftmark}}
\setstretch{1.3}

% Footnotes as the 1928 printing marks them: 1), 2), ... numbered afresh
% on each printed page.  Always written \footnote[N]{...} with the
% printed number.
\renewcommand{\thefootnote}{\arabic{footnote})}

% Editorial: a stretch of the book not yet transcribed.  Remove the \gap
% line together with its marker when the stretch is transcribed
% (splice.py does this).
\newcommand{\gap}[1]{\par\begin{center}\small[#1 not yet transcribed]\end{center}\par}
% A white line in the print, and a short centred rule.
\newcommand{\whiteline}{\par\vspace{\baselineskip}\par}
\newcommand{\secrule}{\par\begin{center}\rule{2cm}{0.4pt}\end{center}\par}
% Editorial: letters supplied where the SCAN has lost them (not the print).
\newcommand{\supplied}[1]{\textlangle#1\textrangle}

\usepackage{hyperref}
\hypersetup{
  pdftitle={Erindringer},
  pdfauthor={Harald Høffding},
  hidelinks
}

\usepackage{marginnote}
\newcommand{\opage}[1]{\marginnote{\footnotesize\textit{[#1]}}}
\newcommand{\apage}[1]{\marginnote{\footnotesize\textit{[#1]}}}

\begin{document}

\frontmatter
% TITLE PAGE (PDF 6).  Read at the image; the letterspaced capitals are
% given as \emph.  The fleuron and the thick-and-thin rule are shown
% schematically.
\begin{titlepage}
\centering
\vspace*{2cm}
{\small\emph{HARALD HØFFDING}}\par
\vspace{1.5cm}
{\LARGE\emph{ERINDRINGER}}\par
\vspace{1cm}
{\large\textasteriskcentered}\par
\vfill
\rule{0.6\textwidth}{1.2pt}\par\vspace{1pt}\rule{0.6\textwidth}{0.3pt}\par
\medskip
{\small\emph{GYLDENDALSKE BOGHANDEL - NORDISK}}\par
{\small\emph{FORLAG - KØBENHAVN - MCMXXVIII}}\par
\end{titlepage}

\chapter*{Forord}
\addcontentsline{toc}{chapter}{Forord}
% Forord, PDF 8, unnumbered; p. [6] by sequence.
\gap{Forord (p.~[6])}
% [text to be added: pp. 6--6]

\mainmatter

% ---- Chapter I.  Indhold: "I. 1843—1861 .... 7".  pp. 7--37 (PDF 9--39).
\chapter*{I}
\addcontentsline{toc}{chapter}{I. 1843--1861}
\markboth{I. 1843--1861}{I. 1843--1861}
\gap{Chapter I, pp.~7--37,}
% [text to be added: pp. 7--37]

% ---- Chapter II.  Indhold: "II. 1861—1874 .... 38".  pp. 38--93 (PDF 40--95).
\chapter*{II}
\addcontentsline{toc}{chapter}{II. 1861--1874}
\markboth{II. 1861--1874}{II. 1861--1874}
% Chapter II is divided for the batches of 23 Sep 2026.  A \gap line
% stands above each stretch not yet transcribed.
\gap{p.~38}
% [text to be added: pp. 38--38]
% p. 39 opens mid-sentence: "de bestaa sammen!"; p. 38 not yet transcribed
% --- p. 39 ---
\opage{39}de bestaa sammen! Maatte der dog ikke tilsidst her kunne
gives et eneste Udtryk for Personlighedens Liv? ---

Det var Rasmus Nielsen og Sibbern, der bleve mine første
filosofiske Lærere. Rasmus Nielsen stod dengang paa sit
Højdepunkt. Han var endnu søgende, betagen af Kierkegaard
og optagen af matematiske og naturvidenskabelige
Studier, stræbende at finde en Vej, ad hvilken disse forskellige
Tankeverdener kunde forenes. Hans propædeutiske
Forelæsninger bestode i mit Aar dels i en Gennemgang af
Kierkegaards „Stadier“, paa Grundlag af et lille Udvalg, han
havde ladet trykke, dels i et kort erkendelsesteoretisk Omrids.
Han talte med Liv og Kraft, ofte aandfuldt og vittigt,
og jeg skylder disse Forelæsninger overordentligt meget.
De satte en Mængde Tanker i Bevægelse hos mig, som
hidtil ikke havde rørt sig. Det gik ham dog med den Forelæsning,
som det saa ofte gik ham med hans Bøger; Udkastet
var for stort og omfattende, eller ogsaa udarbejdede
han de første Dele af det saa vidtløftigt, at der ikke blev
Tid nok til de senere Dele. Af „Stadierne“ fik han saaledes
kun gennemgaaet det første, det „æstetiske“. Ved enhver
mulig Lejlighed polemiserede han imod Teologien som Videnskab,
idet han skarpt skelnede mellem Tro og Teologi.
Denne Polemik vakte til at begynde med min Modsigelse.
Jeg var mødt med store Forventninger om, hvad det teologiske
Studium skulde bringe; jeg haabede paa, at den Isolation,
i hvilken det Religiøse hidtil havde staaet i mit Sjæleliv,
ved Teologiens Hjælp vilde vige for en Harmoni mellem
de forskellige Sfærer. Jeg haabede at blive ført ind i
en Forening af Tankeliv og Følelsesliv, saaledes at de supplerede
hinanden; jeg haabede, at Atene og Madonna kunde
forenes til Et. I en Dagbog, jeg holdt i mine første Studenteraar,
gav jeg da ogsaa ondt af mig om dem, der ikke
vilde anerkende, at der gives en kristelig Videnskab. --- Det
skulde snart blive anderledes.

Medens Rasmus Nielsen læste over filosofisk Propædeutik,
læste Sibbern over Psykologi og Logik. (Det filosofiske
% --- p. 40 ---
\opage{40}Kursus lagde dengang Beslag paa ti Timer om Ugen. Det
var maaske for meget. Men man er gaaet for vidt ved at
begrænse det til fire Timer, selv om der nu ved Brug af
trykte Grundlag kan udrettes ikke saa lidt i den Tid. Efter
min Mening vilde sex Timer være passende). Naar Propædeutik
skiltes fra Psykologi og Logik, var det, fordi den
nærmest skulde give Erkendelsesteori eller almindelig Videnskabslære.
Rasmus Nielsen gav forresten meget Andet
foruden Erkendelsesteori under Rubriken Propædeutik; saaledes
i mit Aar Oversigten over de Kierkegaardske Stadier.
Der var en tredie Professor, Hans Brøchner; han dublerede
skiftevis Sibbern og Nielsen; og i mit Aar læste han over
Sibberns Fag. Jeg havde valgt Sibbern, da jeg gerne vilde
høre ham, inden han gik af, og da jeg ventede mig meget
af Forfatteren til „Gabrielis“. Han var dengang 77 Aar.
Det blev en Skuffelse. Der kom vel af og til geniale Glimt,
smukke Tanker og kostelige Bemærkninger; men hans
Fremstilling var mat og mekanisk, og hans sløvede Sanser
lod ham hverken høre, hvorledes han selv raslede med sine
Nøgler, mens han talte, ej heller mærke Tilhørernes Uro
og Snakken. Hans Naivitet og Særhed kunde være kostelige,
saaledes naar han nævnede sin „Gabrielis“ blandt
„gode poetiske Værker“, eller naar han udførligt forklarede
os, hvorfor Frue Kirkes fire Ure ikke gik ens. Men naar
han, fordi han havde det for varmt paa Katedret, først tog
sin Kalot af, og naar det ikke hjalp, sin Paryk, kunde man
ikke forlange, at unge Mennesker skulde tage det med Andagt.
Han havde overlevet sig som Lærer, skønt hans indre
Liv endnu rørte sig med Friskhed. Hans Bøger var drøje
at lære, tilmed da han helst vilde have sine egne Ord igen.
--- Først langt senere, da jeg selv var bleven selvstændig i
Filosofien, har jeg ret lært at paaskønne ham, og det har
været mig en stor Glæde at fremdrage hans Minde og karakterisere
% footnote 1) on p. 40
ham\footnote[1]{Se „Mindre Arbejder“, Første Række, og „Danske Filosofer“.}, samt bidrage til, at hans Tanker ogsaa
bleve bekendte i Udlandet.

% --- p. 41 ---
\opage{41}Jeg begyndte i mit første Studenteraar ogsaa paa mit teologiske
Studium, især Exegese og Kirkehistorie. Det exegetiske
Studium er vel det, jeg har havt mest Gavn af
indenfor Teologien. Det vænnede til Analyse af, hvad man
læste, til nøjagtig og indgaaende Forstaaelse, og det gav
Anledning til interessante Sammenligninger mellem de nytestamentlige
Forfattere. Og det at blive hjemme i et saa
stort Aandsværk som det nye Testamente, at kunne vandre
gennem det paa Kryds og Tvers og føle Tankernes Sammenhæng
med den sproglige Form, var allerede et stort
Udbytte, og det er især dette, der har gjort, at jeg aldrig
har fortrudt at have studeret Teologi, saa langt jeg end er
kommen bort fra teologisk Tankegang. Ganske vist, de vidtløftige
tyske Kommentarer, vi studerede, vare ikke lystelige
at faa Forstand af. Man gennempløjede en hel Række af
Fortolkninger, der i Forfatterens Øjne vare urigtige, og
naaede saa tilsidst Forfatterens egen Tydning, der naturligvis
var den højeste Visdom. Der var et Pedanteri heri,
som kunde give Afsmag for Studiet. Og naar man ret besindede
sig paa, at de Ord, som var Genstand for den megen
Lærdom, gik ud paa at stille dem, der læste dem, overfor
Livets største Spørgsmaal, Spørgsmaalet om aandeligt
% PRINTER'S ERROR: "andelig" so printed (sense: aandelig)
Liv eller andelig Død, saa kunde man ofte harmes over
den hele Behandlingsmaade, tilmed naar man lige havde
% PRINTER'S ERROR: "Kierkegaard" so printed, without the genitive s (sense: Kierkegaards)
læst en Bog som Kierkegaard „Til Selvprøvelse“. Jeg har
flere Gange kastet Kommentarerne henad Gulvet --- men jeg
gik siden hen og tog dem op. Jeg vilde dog fuldføre, hvad
jeg havde begyndt. ---

Under alt dette levede jeg dog foreløbigt som en glad
Student og havde mange fornøjelige Aftener i den gamle
Studenterforeningsbygning i Boldhusgade trods dens Uappetitlighed.
Lørdag Aften var der ofte „Sold paa Salen“.
Man sad med sin Pibe ved Punchebollen; Visebogen med
Plougs og Hostrups Sange sloges op, og mellem Sangene
blev der holdt mere eller mindre patetiske eller vittige Taler.
Diskussioner holdtes ikke, uden paa Generalforsamlin%
% --- p. 42 ---
\opage{42}ger, hvor der kæmpedes for at styrte eller hævde Senioratet.
Ogsaa om Formiddagen kunde man vente at træffe
Kammerater deroppe, og det gode Bibliotek havde jeg megen
Nytte af.

Da der ikke var saa mange Studenter som nu, og
da de havde flere Forelæsninger fælles, var der megen Anledning
til at gøre Bekendtskaber. Jeg kom hurtigt til at
kende Studenterne fra Aalborg Skole, af hvilke flere bleve
mine gode Venner. Der var ikke gaaet mange Dage efter
Forelæsningens Begyndelse, før jeg kom paa fortrolig Fod
% footnote 1) on p. 42
med Andreas Faartoft\footnote[1]{Senere Distriktslæge i Nykøbing paa Mors.}, ret en Type paa en ung Student
fra den Tid, aaben og glad og med stort Mod paa Livet.
Han havde en smuk Sangstemme, som ofte var til Glæde i
mit Hjem, hvor han snart blev en velset Gæst. En anden
% footnote 2) on p. 42
Aalborgenser, Søren Svendsen\footnote[2]{Senere Forstander for Navigationsskolen paa Bogø.}, gjorde jeg ogsaa hurtigt
Bekendtskab med, og han og hans Violoncel bleve ogsaa
snart velsete Gæster i mit Hjem. Mine to unge Kammerater
betragtedes af mine Forældre og Søskende som hørende
til hjemme hos os, naar der skulde være en fornøjelig og
hyggelig Aften. Jeg skylder dem meget, ikke mindst i de
følgende Aar, hvor jeg ikke var noget oplivende Element
i mit Hjem, og hvor jeg dog selv trængte til Oplivelse.
Faartoft blev senere en anset Læge i sin Fødeby, Nykøbing
paa Mors, og Svendsen, der havde studeret Fysik, blev
Navigationsskolebestyrer paa Bogø. --- En tredie Aalborgenser,
% footnote 3) on p. 42
% SETTLED 23 Sep 2026 (2nd reader undecidable; 3rd reader: a point, PROBABLE -- a point smeared down to the right;
% this fount's commas curl down to the left and hang ~10 px lower).  Transcriber's note: "Kierkegaard." vs "Kierkegaard," at the end of note 3 -- the mark is battered (as is the whole line's foot); it sits on the baseline and does not descend like a comma
Poul Kierkegaard\footnote[3]{Søn af Biskop P. Chr. Kierkegaard.}, kom jeg noget senere i nærmere
Forhold til og vil i en anden Sammenhæng fortælle om
ham og hans Skæbne. ---

Samme Dag, jeg havde været oppe til Filosofikum, ankom
norske og svenske Studenter. Det var det sidste skandinaviske
Studentermøde, der blev holdt i København (1862).
Der var stor Begejstring i Byen og mellem Studenterne,
og det blev glade og fornøjelige Dage, der fulgte, Dage, jeg
% --- p. 43 ---
\opage{43}endnu tænker paa som festlige Minder, skønt der snart
faldt mørke Skygger hen over de Forhaabninger, som dannede
Grundlaget for Feststemningen. Jeg mindes særligt
Turen i Nordsjælland, da vi, efter at have overnattet i Hillerød,
næste Morgen vandrede den herlige Vej fra Hillerød
til Fredensborg. Her var der Lejlighed til Samtale og
Tankeudvexling med de skandinaviske Brødre, og man
kunde faa Oplysning om Studieforhold ved de andre nordiske
Universiteter. Helt indtagen var jeg allerede den første
Dag i en ung Svensker, hvis Navn var Svartling, og
som jeg i flere Aar korresponderede med. Mange Aar efter
hørte jeg ham beskrive som en anseelig Provst. Paa Fredensborg
defilerede vi gennem Slottets Kuppelsal, fra hvis
Galleri Frederik den Syvende, omgivet af Grevinde Danner,
Ministre og Hofmænd, modtog vor Hyldest. Vi marcherede
ned i Nordmandsdalen, hvor en Sexa var anrettet.
Men neppe havde vi begyndt at tage for os af Retterne, før
en voldsom Tordenbyge brød løs. Vi dukkede ned under
Bordene, medtagende hvad vi kunde faa fat i af Spise- og
Drikkevarer, især af de sidste. Intet Under, at Stemningen
var høj, da Solen igen kom frem, og vi dukkede op igen.
Lidt efter kom Kongen, og paa sin gemytlige Maade sagde
han, at han haabede, vi vare bleven ligesaa vaade indvendigt
som udvendigt. Turen fortsattes derefter til Helsingør.

Den Sommer boede min Familie i et Hus ved Krogerup
Allé, hvorfra der var en herlig Udsigt over Markerne til
Sundet og til Helsingborg. Det er en af de bedste Sommerferier,
% footnote 1) on p. 43
jeg har havt\footnote[1]{Jeg ser af mine Breve til Hans Skovgaard fra Sommerferien 1862,
at jeg har beskæftiget mig meget med Aristofanes (i dansk Oversættelse),
særlig for at faa Klarhed om Forholdet mellem ham og Sokrates. Om
Sokrates drejede min Interesse sig allerede fra Skolen af. Desuden har jeg
læst Platon’s „Gorgias“ (dels paa Græsk, dels paa Dansk), Rötscher’s „Aristophanes
und sein Zeitalter“, Goethe’s „Aus meinem Leben“ og Kierkegaards
religiøse Taler. Det var Forbindelsen af disse Beskæftigelser med
Udflugter i den herlige Egn, der gjorde Ferierne i Humlebæk uforglemmelige.}. Vi strejfede om i den herlige Egn,
snart til Hest og snart til Fods. Forpagteren paa Krogerup
% --- p. 44 ---
\opage{44}havde to islandske Heste, han overlod min yngre Broder
og mig, og vi kunde saaledes komme langt omkring paa
en fornøjelig Maade. --- En Dag kørte vi til Danstrup Hegn
ved Fredensborg og saa her Frederik den Syvende i Spidsen
for sin Hestgarde modtage Karl den Femtende af Sverige,
der kom kørende fra Helsingør og nu steg til Hest,
hvorpaa begge Konger i Spidsen for et glimrende Følge
rede til Fredensborg. Dette Kongemøde mentes, ligesom
det forudgaaende Studentermøde, at skulle indvarsle en ny
Tid for Norden. Ak, det var Sæbebobler, og der kom tunge
og mørke Aar for vore ungdommelige Idealer.

\whiteline

Der kom nu ogsaa en tung og mørk Tid i mit indre
Liv. Jeg fordybede mig i Kierkegaards Skrifter, og efter
en haard Kamp maatte jeg give ham Ret i hans Opfattelse
af Kristendommen. Denne stod ikke længere som den
milde, trøstende Sandhedsvej, der i sig kunde optage alle
menneskelige Livsveje --- Familieliv, Kunst, Videnskab,
Samfundsliv --- og føre dem hen imod Fuldendelsen. Idyllen
% PRINTER'S ERROR: "sprængtes.Hinsidighedens" so printed, no space after the full stop (sense: sprængtes. Hinsidighedens)
sprængtes.Hinsidighedens Ideal kom til at staa i sin
hele ubetingede Strenghed. Kun Et var i Ordets bogstavelige
Forstand fornødent, det at træde i Forhold til det
store Paradox, Menneskeblivelsens Under, der lod alt
Menneskeligt blegne og kun stillede én Opgave, den at opgive
Alt for at faa Del i de guddommelige Kræfter, som
paa ét eneste Sted i Verden og paa ét eneste Punkt i Tiden
havde gennembrudt Naturens og Menneskelivets Sammenhæng
og ladet Evighedens Lys skinne ind i Mørket. Udsagn
i det nye Testamente, jeg hidtil havde skudt tilside
eller tolket mildt og harmoniserende, kom nu til at staa i
deres hele og absolute Idealitet. Det var ikke Kristendommens
Mening at fuldende, hvad menneskelig Kultur indledede,
men at bryde af og at henstille alle menneskelige
Forhold som ligegyldige, forsaavidt de ikke vare Hindringer
for at naa det store Maal. Det faldt mig som Skæl
% SETTLED 23 Sep 2026 (2nd reader undecidable; 3rd reader: a point, PROBABLE -- the same smear; the comma after
% "Øjnene," on this line is a different shape).  Transcriber's note: "Testamente." vs "Testamente," -- the mark after the last word of p. 44 is battered; it sits on the baseline and does not descend like the comma of "Øjnene," on the same line, and p. 45 opens with a capital
fra Øjnene, da jeg nu igen tog fat paa det nye Testamente.
% --- p. 45 ---
\opage{45}Og det stod for mig, at det, som dér forkyndtes, var det
Højeste. Jeg kan ikke henvise til noget bestemt Omvendelsesøjeblik.
Det var efter haard Modstand, at jeg opgav
min første Ungdoms harmoniske og idylliske Tro. At
Kristendommen var og vilde være Trøst, den evige Trøst,
var sandt; men denne Trøst var knyttet til én Betingelse:
fuld Hengivelse til det ene Fornødne, fuld Opgivelse af
Livets menneskelige Formaal. Jeg opdagede egentligt først
nu, at hverken Kunst eller Videnskab eller borgerlig Virken
omtales i det nye Testamente, og at der advares i
strenge Ord mod Ægteskab. Det var den kristelige Etik,
der nu viste sig for mig i et helt nyt Lys. Med Dogmatiken
var der foreløbigt Intet ivejen.

Hvorledes skulde nu Livet føres, efterat Idealet saaledes
havde vist sig at ligge højere, end jeg havde tænkt? Den
Tid, som nu kom, bragte mig vel Øjeblikke af inderlig Begejstring,
men især bragte den mig Tvivl, Uro og Hjemløshed.
Jeg vidste ikke mine levende Raad. Mesteren, Profeten
var jo gaaet bort, efterat han havde sagt, at det nye
Testamentes Kristendom ikke var til, om den overhovedet
havde været til ud over den allerførste Tid. Det var forgæves
at spekulere over, hvorledes han vilde have stillet
sig, hvis han ikke saa tidligt var bleven bortreven. Og desuden
havde han jo selv lært mig, at hver Enkelt maa løse
Livsproblemet ad sin egen Vej, --- at man i Grunden ikke
kan lære af Andre.

Og nu mit teologiske Studium! Jeg arbejdede med Iver i
de forskellige teologiske Fag, men fortsatte ogsaa de filosofiske
Studier, hørte stadigt Rasmus Nielsens Aftenforelæsninger,
og nu ogsaa Brøchners over Filosofiens Historie.
Jeg deltog i et Privatissimum, Brøchner holdt over
Kant’s Kritik der Urtheilskraft. Men fra rent filosofisk
Side blev det stedse mere hævdet, at en videnskabelig Harmoni
mellem Tro og Viden var umulig, --- at der i Grunden
ingen teologisk Videnskab gaves. Saa vandrede jeg
da frem og tilbage mellem Teologi og Filosofi --- ogsaa
% p. 46 opens mid-sentence, not mid-word: p. 45 ends "ogsaa", p. 46 begins
% "bogstaveligt talt"; the paragraph continues from p. 45 (no indent)
% --- p. 46 ---
\opage{46}bogstaveligt talt, idet de havde deres Auditorier hver i sin
Ende af den lange Universitetsgang ovenpaa. Der gjorde
sig en stadig intellektuel Splid gældende hos mig, foruden
den personlige, som fremkaldtes ved den nye Opfattelse
af Kristendommen.

Det var de tungeste og mørkeste Aar i mit Liv. Jeg lukkede
mig inde med mine Tvivl og blev uomgængelig og tillukket
overfor mine Nærmeste. Hvor meget jeg end holdt
af dem, var jeg ikke noget godt Familiemedlem. Livet paa
mit Kammer og Livet i Dagligstuen kom til at staa for
mig som to Verdener, og jeg formaaede kun altfor slet at
gøre Overgangen mellem dem. Den lyse og harmoniske
Maade, hvorpaa de dernede toge Alt, opirrede mig. Der
kom netop i de Aar forøget Velstand i Huset, og man nød
paa sin Vis, hvad den nu kunde aabne Adgang til. Jeg forholdt
mig afvisende og kold. Mange Synder har jeg paa
min Samvittighed fra disse Aar, og først senere hen har
jeg forstaaet, i hvor høj Grad man har maattet bære over
med mig. Mine Forældre vare, som jeg senere har erfaret,
bekymrede for mig og spurgte mine Venner tilraads.

Under alt dette virkede ogsaa et Hang til Hypokondri,
% PRINTER'S ERROR: "mine legemlige Konstitution" so printed (sense: min)
som hænger sammen med mine legemlige Konstitution og
har fulgt mig Livet igennem, helt op i de gamle Aar. Kun
mine Nærmeste have havt Lejlighed til at mærke dette
Hang, som det overfor Fremmede i Regelen falder mig
let at trænge tilbage. I min Bog „Den store Humor“ er det
Stykke, der handler om Hypokondri (acedia), et af de Afsnit,
der mest ere byggede paa Selverfaring. Tunghed og
Mørke kunde uden paaviselig Aarsag brede sig over Sindet
og give Alt et mat og taaget Udseende. I lang Tid kom
der ofte en Fornemmelse op, som om jeg hang i et Haar,
der hvad Øjeblik det skulde være kunde briste. Utilfreds
med mig selv og over, at jeg intet udrettede i Verden, syntes
jeg ofte, at Alt var haabløst for mig. Dette Hang
maatte naturligvis saa meget des lettere faa Overvægt, som
Sindet stod tvivlende og uroligt overfor de største Spørgs%
% --- p. 47 ---
\opage{47}maal. Hvorledes Styrkeforholdet var mellem de to Strømninger,
den ene, der kom fra den legemlige Tyngde, og den
anden, der kom fra de intellektuelle og de personlige Problemers
Tyngde, kan jeg naturligvis ikke udrede, og kunde
det endnu mindre dengang.

Men det forholdt sig naturligvis ikke saaledes, at jeg stadigt
gik omkring som en Fortabt. Mine Studier optog mig
meget. Foruden mit Fagstudium fik den Interesse for Hebraisk,
der allerede var vakt i Skolen, mig til at læse Arabisk
hos Professor van Mehren. Jeg kom saa vidt, at jeg
fik læst et Par Fortællinger af „Tusind og een Nat“ paa
Grundsproget, og jeg tænkte endog en Stund paa at blive
Orientalist, hvad der vel vilde have virket afledende paa min
Betagethed af filosofiske og religiøse Problemer. Jeg læste
endnu adskillig Filosofi, særligt forfulgte jeg Udviklingen
fra Kant til Hegel, en Periode, der især paa den Tid stod
med særegen Glans. Ikke blot havde jo Kierkegaard været
stærkt paavirket af Hegel, men baade Rasmus Nielsen og
Brøchner betegnede Nuancer af Hegelianismen. Jeg læste
ogsaa nogle Dialoger af Platon paa Grundsproget. Kort
sagt, min Tid var godt besat.

Systematisk Teologi læste jeg sammen med Hans Skovgaard,
i Regelen ude i hans Hjem paa Valby Bakke. Vor
Læsning strakte sig ofte langt ud paa Natten, og jeg havde
da en Vandring for mig i Bælgmørke fra Bakken ned gennem
den snævre Smøge, der førte fra Langgade til Rahbeks
% "ligened" so printed, closed up as one word (read at 600 dpi; sense: lige ned)
Allé, og ligened til Vesterbrogade, hvor Civilisationen
med sine Lygter begyndte. Vejen var ikke ufarlig, da der
ofte forekom Overfald udenfor Byen, og da der i Landstedet
lige for Pilealleen (nu Ny Carlsbergs Hovedbygning)
gik nogle store Hunde løse, som fore ud mod den natlige
Vandrer, uden at han kunde se dem. Men det var en forfriskende
Vandring, jeg havde derud om Eftermiddagen.
Jeg erindrer en Sommereftermiddag, da jeg som sædvanligt
kom for at læse et aftalt Afsnit igennem. Jeg traf Skovgaards
næstyngste Søster i Havestuen, og hun --- som altid
% --- p. 48 ---
\opage{48}var elskværdig --- viste mig den Dag en ganske særlig Venlighed
og fik mig indviklet i en lang og livlig Samtale. Bagefter
kom det for Dagen, at hun agerede efter Aftale med
Broderen, der ikke var bleven færdig med sit Pensum. Hun
maatte ofte høre for den Komedie, hun havde spillet. Hun
var en sund og frejdig ung Pige, og jeg tænker tit paa
hende, naar jeg gaar forbi det gamle Bjerregaard. For
nogle Aar siden havde jeg hendes Datterdatter, som i høj
Grad lignede hende, blandt mine Tilhørere ved Universitetet.
--- Det var baade fornøjeligt og lærerigt at læse sammen
med Hans Skovgaard. Han var meget interesseret,
men tillige meget skeptisk og oppositionel. Jeg var dengang
i Grunden ivrigere Teolog end han; senere er det blevet
omvendt. Han var især ikke saa ivrig som jeg til at finde
Mening i Professor I. A. Bornemann’s ofte dunkle Ord.
Det er efter min Erfaring noget, der ikke noksom kan tilraades
unge Studerende, at læse sammen med Kammerater.
Derved udvikles Selvtænkning og Kritik mere end ved
Læsning med Manuduktør, der, naar Manuduktøren ikke
er særligt pædagogisk anlagt, let faar en mekanisk Karakter.

Foruden Valby var der et andet Sted, til hvilket jeg ofte
vandrede ud, og hvor jeg fandt Hvile og Forstaaelse. En
gammel Onkel af mig boede i et Landsted i Frederiksberg
Smallegade, og hans Hus blev bestyret af en fjern Slægtning
af mine Forældre, Sofie Pape, Søster til den unge
% footnote 1) on p. 48
Pige, jeg havde gæstet mod Slutningen af min jyske Rejse\footnote[1]{De to Søstre var Døtre af Proprietær H. C. Pape til Gaarden Heiselt
i Vendsyssel. En Tid havde han tillige en Forpagtning i Nærheden af Viborg.
(Se ovenfor p. 36).}.
I hende fandt jeg en ældre Søster, som jeg ofte havde ønsket
mig. Hun forstod min Interesse for Kierkegaard, idet
hun dog væsentligt holdt sig til hans religiøse Taler. Og
nu kom i Vinteren 1863---64 hendes yngre Søster Emma i
Besøg hos hende. Vi tre kom til at danne en lille Klike.
Men det blev snart for den yngres Skyld, jeg kom der. Vi
læste sammen og vi spadserede sammen. Og jeg mente nu,
% --- p. 49 ---
\opage{49}at jeg havde faaet en jævnaldrende Søster. At det var en Følelse
af anden Art, der rørte sig hos mig, blev mig først klart
efter hendes Bortrejse. Jeg trængte foreløbigt Alt, hvad der
rørte sig i den Retning, tilbage under Indflydelse af den
strenge religiøse Tugt, jeg havde givet mig ind under. Og
dog var, langt mere end jeg vilde være ved overfor mig
selv, Tanken om hende i disse mørke Aar et Lysskær, som
holdt Mod og Haab oppe. Der laa desuden her Spiren til
noget, der senere virkede til at bringe mig i et levende
Forhold til Livet igen.

Vinteren 63---64 var mørk og tung for mit Fædreland.
Med Spænding fulgte man Begivenhederne, som de udfoldede
sig efter Frederik den Syvendes Død, den truende
Forbundseksekution, Oprøret i Holsten og Preussens og
Østerrigs Indgriben. Saa i Februar Kampene ved Dannevirke
og Hærens pludselige Tilbagetog. Sorg og Harme
betog Alle, og efter Dybbøls Fald stod Nederlaget som
uafvendeligt. Det var bitre og mørke Tider. Men jeg skammer
mig ved at tænke paa, hvor meget jeg under alt dette
alligevel gik op i mine indre Brydninger og Grublerier og
lod mig beherske af den Afmagtsfølelse, der var Følgen,
uden at Landets Ulykke kunde vække mig til Handling. I
Forsommeren bestemte jeg mig dog, sammen med min Ven
Skovgaard, til at melde mig til Udskrivning. I nogen Tid
stod det saaledes for mig, at nu skulde Alt lægges tilside,
og det stod saa ogsaa klart for mig, at hvis jeg kom tilbage
fra Krigen, vilde jeg ikke vende tilbage til Teologien.
--- Da der ikke blev antaget flere Soldater, vendte jeg dog
tilbage til Studerekammeret.

Jeg havde ikke opgivet Kristendommen, men jeg gjorde
overfor mig selv den Indrømmelse, som Kierkegaard (i
„Indøvelse i Kristendom“) havde forlangt af Kirkens Repræsentanter.
Jeg anerkendte, at jeg ikke levede under den
Etik, Urkristendommen forudsætter eller medfører. Jeg
erkendte min Afstand fra det urkristelige Ideal, skønt jeg
fastholdt, at det var Idealet. Jeg følte mig som en Katolik,
% --- p. 50 ---
\opage{50}der ikke gaar i Kloster, skønt han indrømmer, at det er
Idealet. Jeg mente, at jeg paa disse Vilkaar havde Ret til
at anse mig selv for Kristen, skønt jeg ikke kunde tænke
paa at blive Præst.

Kirken som Kirke kunde --- det har jeg siden forstaaet
--- ikke gøre den Indrømmelse, Kierkegaard forlangte. Det
vilde have været det samme som at indrømme, at den repræsenterede
en helt anden Religion end den urkristelige,
naar en Religion og dens Etik overhovedet erkendes at høre
saa nøje sammen, at de staa og falde med hinanden. Men
ogsaa den Enkelte, der gør Indrømmelsen, vedkender sig
jo egentligt --- det saa jeg heller ikke klart dengang --- at
han ikke har den Religion, Kirken paastod at repræsentere.
Det betegnede i Virkeligheden et Brud med Kirken i Kraft
af egen Livserfaring og Tænkning, og jeg stillede mig derved
ind under Personlighedsprincippet. Rækkevidden heraf
gik kun sent og langsomt op for mig i Sammenhæng med
min fortsatte Livserfaring og mine fortsatte Studier. Paa
dette Punkt peger Kierkegaard selv, i Kraft af det Princip,
at Subjektiviteten er Sandheden, langt ud over sig selv.
Min Kærlighed til Kierkegaard og min Beundring for ham
har derfor kunnet følge mig længe efter, at jeg var kommen
bort fra alle Dogmer. Det er dog, i Anledning af min
Bog om Kierkegaard, blevet sagt om mig paa Tysk af en
letbenet Teolog, at jeg mangler Kærlighed til Kierkegaard
og Fordybelse i ham. Men Sagen er, at der er Mennesker,
som ikke forstaar en Kærlighed og en Beundring, der
ikke giver sig Luft i sentimentale Udtryk, men i at tage et
stillet Problem op og trofast arbejde med det paa den Maade,
der stemmer med den Enkeltes Personlighed. Kierkegaards
Problem har fulgt mig fra min Ungdom, bestemt
mit Livs Retning, atter og atter ført mig til en inderligere
Opfattelse af mangfoldige Forhold, til en strengere Prøvelse
af mig selv, til en Sky for, hvad ikke har personlig
Sandhed. Og medens Teologer og Halvteologer i Hundredaaret
efter hans Fødsel kalkede Profetens Grav, holdt jeg
% --- p. 51 ---
\opage{51}mig til det store Princip om Subjektiviteten som Sandhed
og søgte at anvende det paa den Maade, der stemte med
min Personlighed og mine Studier.

Rent filosofisk har Studiet af Kierkegaard ledet mig til
en Tanke, mit senere filosofiske Studium førte til at uddybe
og anvende i det Enkelte, den Tanke nemlig, at det aandelige
Livs formale Ejendommelighed ligger i den Enhed,
den Koncentration, hvormed Livsindholdet sammenfattes,
og at Maalestokken for et Aandslivs Værdi ligger i Forholdet
mellem Indholdets Fylde og Koncentrationens
Energi. Denne Tanke, der i Filosofiens Historie er kommen
frem før Kierkegaard, ligger til Grund for Kierkegaards
filosofiske Karakteristik af de forskellige Stadier
eller Livstyper. Dette Punkt stod naturligvis heller ikke
klart for mig. Det gik først op for mig gennem et kritisk
Studium af engelsk Filosofi og gennem Forarbejderne til
min Psykologi.

Der var et andet Synspunkt, fra hvilket jeg i hine Aar
heller ikke endnu kunde bedømme Kierkegaards Opfattelse
af Urkristendommen. Teologerne pleje at betragte den
Tanke, der bærer hele den urkristelige Etik, nemlig Forventningen
om det snarlige Indtræden af „Himmeriges
Rige,“ som en uvæsentlig Illusion hos den ældste Menighed,
eller de bortforklare de Steder, hvor den omtales med
særlig Styrke. Og til min store Forundring saa jeg, da
Kierkegaards efterladte Papirer bleve trykte, at han selv
betragtede Forventningen som en Illusion --- kun at han
vendte Forholdet om: der kunde efter hans Mening ikke
leves paa en Etik som den urkristelige uden hin Illusion:
„Det er i den Grad kvalfuldt at være den sande Kristne,
at det ikke var til at holde ud, hvis man ikke bestandigt
ventede Kristi Genkomst netop nu.“ (Efterladte Papirer.
1849, p. 248). Ved at kalde Forventningen en Illusion stiller
Kierkegaard sig i fuldstændig Modsætning til Urkristendommen.
Fra et historisk Standpunkt er Forventningen
et Faktum og er den stadige positive Forudsætning,
% p. 51 ends mid-sentence at "Forudsætning,"; p. 52 not yet transcribed
\gap{pp.~52--54}
% [text to be added: pp. 52--54]
% p. 55 opens mid-word: "Ved-/kommende" (p. 54 ends "hans Ved-"); p. 54 not yet transcribed
% --- p. 55 ---
\opage{55}kommende. Han stillede et Bundt Cigarer paa sit Bord og
bestemte, at det skulde være de sidste, han røg. Denne
sidste Nydelse opsatte han saa Dag for Dag, og i Løbet
af et Aar var saa Trangen forsvunden.---

Examenstiden gik som Sport, og den 19. Juni 1865
vandrede jeg som candidatus theologiæ ned ad Universitetstrappen.
Om Eftermiddagen havde jeg den tidligere
omtalte Samtale med min Fader om min Fremtid. Mit
Ønske var at slaa ind paa Skolevirksomhed og at fortsætte
mine Studier. Jeg fik, endnu inden jeg var taget
ud til min Familie, som laa paa Landet i Skodsborg, Tilbud
om Timer i „Mariboes Skole“ (ved Frederiksholms Kanal)
i Historie og Dansk, som jeg særligt havde tænkt
mig som mine Undervisningsfag. Ved denne Skole virkede
jeg som Lærer i 17 Aar, lige til jeg blev Professor,
og det er mange gode Minder, jeg har fra den Tid. Foreløbigt
holdt jeg Ferie i Skodsborg, og senere paa Sommeren
foretog jeg en Rejse til Stockholm, hvorfra jeg ad
Kanalvejen vendte hjem over Gøteborg.

Det var svære Problemer, jeg tog med mig fra den
Periode i mit Liv, som nu var sluttet. Jeg havde bevaret
min personlige Tro, med det Forbehold, jeg før har beskrevet.
Men i hvilket Forhold skulde og kunde nu denne
Tro staa til, hvad der ellers optog mig i Livet? Det var
det Hovedspørgsmaal, som fulgte mig. Og til at besvare
det, var filosofisk Tænkning nødvendig. Filosofien, der
hidtil havde været et Bifag for mig, blev nu et Hovedfag.
Ti hint Spørgsmaals Besvarelse krævede en Undersøgelse
af, hvad menneskelig Erkendelse er og formaar, og i hvilket
Forhold den staar til det personlige Liv. Jeg maatte
lære, hvad det menneskelige Aandsliv formaar at naa paa
egen Haand, for saa at kunne se, hvorvidt en Forbindelse
mellem det og den religiøse Tro var mulig, og paa hvilke
Vilkaar. Det var det Problem, Kierkegaardstudierne havde
stillet mig overfor, men ikke havde løst. Jeg maatte forfølge
dette Spørgsmaal videre. Min Fremtid tænkte jeg
% --- p. 56 ---
\opage{56}mig foreløbigt saaledes, at Lærervirksomheden blev min
Gerning udadtil, og at der til den knyttedes de Studier,
den krævede, --- men at jeg saa vilde stræbe at have et
Lønkammer, hvor det, der laa mig personligt paa Sinde,
kunde dyrkes. At min Filosoferen skulde kunne faa nogensomhelst
Betydning for Andre, faldt mig ikke ind.

Og saa begyndte jeg da efter Sommerferien min Lærergerning.
Jeg havde ligefra min Skoletid interesseret mig
for historisk Læsning, og nu planlagde jeg et Studium
af de vigtigste historiske Værker, ligefra Herodot og Thukydid
til Macaulay og Ranke. Jeg tænkte ogsaa et Øjeblik
paa at besvare en Prisopgave, Universitetet havde udsat
om et Emne fra den danske Middelalder; men det blev
ikke til noget. Filosofien pressede paa, og som Kontrast
til den megen Teologi, jeg havde gennemgaaet, følte jeg
Trang til at fordybe mig i græsk Filosofi. Som jeg nogle
Aar senere skrev i mit Doktorvita: „Den antike Oldtid
traadte frem for mig i et friskere og klarere Lys netop
ved Modsætningen til den skolastiske Periode, jeg forekom
mig at have gennemlevet; jeg glædede mig ret ved
at læse dens Historieskrivere, dens Digtere, og snart især
ogsaa dens Filosofer. Jeg vandt derved en dyb Overbevisning
om Videnskabens og det humane Livs selvstændige
Betydning.“ Ogsaa her virkede stadigt Kierkegaards Vejledning.
Han havde klart fremholdt den sokratisk-platoniske
Opfattelse af Tanke og Liv som en Aandsmagt, man
maatte gøre sig Rede for, inden man med fuld Ret og med
klar Bevidsthed kunde komme til et Stade, der virkeligt
var højere. Jeg læste nu de græske filosofiske Værker i
Sammenhæng, og det blev i de følgende Aar mit Hovedstudium.
I min Lærervirksomhed ombyttede jeg nogle Aar
senere Historie og Dansk med Latin (i Mellemklasserne)
og Græsk (i de øverste Klasser), og der kom saaledes en
vis Sammenhæng mellem mit Liv paa Skolen og i Studerekammeret.
% "Maribos Skole" so printed here; p. 55 prints "Mariboes Skole".
Foruden i „Maribos Skole“ underviste jeg ogsaa
i „Lyceum“ og en kort Tid i „Bohrs Skole“. Jeg holdt
% --- p. 57 ---
\opage{57}meget af at undervise, især i de øverste Klasser, og har
her vundet Venner, der ogsaa efter deres Skoleaar have
vist mig trofast Hengivenhed.

\whiteline

I Efteraaret 1865 kom den unge Pige, Emma Pape, der
for et Aar siden var rejst bort som min „Søster“, igen til
Byen for at tilbringe Vinteren der. Strax da jeg hørte,
at hun skulde komme, stod det klart for mig, at der nu
maatte komme en Afgørelse, da hun i Virkeligheden var
andet og mere for mig end en Søster. Jeg var bange for
en Afgørelse, ikke blot som enhver ung Mand vel er det,
naar han ser et Enten-Eller nærme sig, men ogsaa af
Grunde, der hang sammen med mit religiøse Stade. Jeg
havde for mit Vedkommende gjort den „Indrømmelse“,
Kierkegaard forlangte, og som hindrede mig i at blive
Præst. Men denne Indrømmelse udelukkede ikke Uro
overfor det Spørgsmaal, hvorvidt fuld Deltagelse i det
naturlige Menneskelivs Interesser og Opgaver var forenelig
med en strengt kristelig Livsførelse. De stærke Steder
i det nye Testamente angaaende Ægteskab, og mod Ægteskab,
kunde jeg ikke komme fra. Og der stod nu, fremkaldt
ved det Enten-Eller, der nærmede sig, en sidste
% "Korinther-/brev" hyphenated at the line end; joined as one word.
Kamp i mit Indre. Det syvende Kapitel i første Korintherbrev
mødte op i hele sin Vælde, og Kierkegaards „Øjeblikke“
laa paa min Pult. Jeg var ikke „søgende“ i samme
Forstand, som Kierkegaards „cand. theol. Ludvig From“,
men jeg søgte rent personligt at klare mig, hvad Konsekvensen
var af, hvad jeg holdt fast paa som det Højeste.
Det var det mest afgørende Vendepunkt i mit Liv, jeg
stod ved. Den naturlige, menneskelige Trang sejrede, og
dermed blev i Virkeligheden Grunden lagt til hele min
senere Maade at tage Livet paa. Som jeg tidligere havde
slængt de teologiske Kommentarer hen ad Gulvet, saaledes
slængte jeg nu Kierkegaards „Øjeblikke“. Ogsaa dem tog
jeg op igen, men fra nu af gik min Beskæftigelse med
Kierkegaard ud paa at finde de psykologiske og historiske
% --- p. 58 ---
\opage{58}Forudsætninger for hans Optræden og paa at hente rent
menneskelig Styrkelse for mit Liv ved at drage Konsekvensen
af, at Subjektiviteten er Sandheden. ---

Og saa var det en Eftermiddag i Efteraaret, medens
jeg stod bøjet over min Thukydid, at det kom til at staa
for mig: nu eller aldrig. Uden at give mig Tid til at
lukke Bogen ilede jeg ud til min Onkels Landsted paa
Frederiksberg, og her blev i den stjerneklare Oktoberaften
den Pagt sluttet, som indledte den korte huslige
% The square brackets are printed (the author's own; cf. "[Senere
% Anmærkning]" on p. 324).
Lykke, der skulde falde i min Lod [i mine unge
Dage].

Min Forlovede blev Vinteren over i København. I de
følgende fem Aar, der gik, inden vi kunde gifte os, var
hun omtrent hvert andet Aar i København, nu i mine
Forældres Hus; de andre Aar tilbragte jeg mine Ferier
enten i hendes Onkels, Provst Wulffs Præstegaard i
Gunderup ved Aalborg, eller hos hendes Forældre paa
deres Gaard i Vendsyssel.

Først ved, mange Aar efter hendes Død, igen at læse
hendes Breve fra Forlovelsestiden, er det ret gaaet op for
mig, hvor uheldigt hun var stillet i disse Aar, naar hun
ikke var i København. Hendes Onkel og Tante i Gunderup
var brave og velmenende Mennesker. Men Provstinden
var af den gamle Skole. Hele Dagen optoges af
Husgerning og Syarbejde. Spadsereture vare Luxus, og
ligesaa det at læse eller skrive Breve. Hjemme hos sine
Forældre havde Emma mere Frihed til at skrive, læse og
% "spasere" so printed (cf. "Spadsereture" two lines above).
spasere; men hendes Moder var meget vanskelig at omgaas,
og der var mange pinlige Tider. Flere Gange ønskede
Emma at tage en Plads som ung Pige i et fremmed Hus,
men jeg var altid imod det. Jeg har dengang ikke rigtigt
gjort mig levende, hvor vanskelig hendes Stilling var.
Gode Tider havde vi dog, dels i København, dels i Sommerferier
paa hendes Forældres Gaard i en karakteristisk
vendsysselsk Egn, med vid Udsigt over de kraftige og
bakkede Omgivelser.

% --- p. 59 ---
\opage{59}I Vinteren 1865---66 førtes der i Pressen en Diskussion
om Rasmus Nielsens Lære om Tro og Viden, som to helt
forskellige Principer, der netop paa Grund af deres Forskellighed
skulde kunne være praktisk og personligt forenelige.
Jeg var ført til at dele denne Opfattelse gennem
min egen personlige Udvikling i religiøs Henseende, saavel
som gennem mine filosofiske Studier. Dette Sidste
trænger til en kort Forklaring.

% "dér" with an acute accent, so printed (read at 600 dpi).
Herhjemme blev endnu omkring 1860 Hegel’s Filosofi betragtet
som et Højdepunkt af Tænkning. Der var dér gjort
et storartet Forsøg paa at vise, at ligesom i den menneskelige
Aand den ene Tanke kun forstaas ved dens Forhold til
andre Tanker, saaledes er i Naturen, i Historien, i Erfaringens
Verden idethele, intet Fænomen, ingen Begivenhed
isoleret, men Alt paa den inderligste Maade forbundet indbyrdes.
Som den ene Tanke fremgaar af den anden, saaledes
det ene Erfaringsfænomen af det andet. Og Hegel
mente, at det var muligt at udvikle den indbyrdes Sammenhæng
mellem Tankerne saaledes, at den tillige kom
til at udtrykke Sammenhængen mellem Tingene. Deraf
udsprang hans store Tillid til den rene Spekulation. Den
menneskelige Tanke kunde ikke være noget Enestaaende
i Verden; dens Love maatte være Verdenslove.

Herhjemme havde denne Lære været Genstand for indgaaende
Kritik af Sibbern og Kierkegaard. De Antydninger,
som disse Tænkere gav af en mere kritisk og erfaringsmæssig
Filosofi, bleve ikke paaagtede. Foreløbigt
stod det nu saaledes, at baade Videnskaben og Religionen,
hver for sig, gjorde Fordring paa at give en absolut Løsning
af Tilværelsens Gaader. Her stod det for Hegel saaledes,
at de religiøse Forestillinger i Billedets Form indeholdt
det Samme, som Filosofien fremstillede i Tankens
rene Form. I Tillid til denne Opfattelse mente mange
Hegelianere at være i god Overensstemmelse med Kirkelæren,
% PRINTER'S ERROR: "Rasmus Nielsens hørte" so printed (sense: Nielsen); read at 600 dpi.
og Rasmus Nielsens hørte en Tidlang til deres Tal.
Men der kom et stort Omslag i hans Opfattelse under Ind%
% --- p. 60 ---
\opage{60}flydelse af Kierkegaards Skrifter. Han brød med sine tidligere
Meningsfæller, navnlig Martensen, og hævdede nu
Religionens absolute Selvstændighed overfor Videnskaben.
De kom nu til at staa som to forskellige Verdener, og
Mennesket maatte leve i dem begge. Denne Opfattelse
blev angreben fra to Sider. Teologien krævede stadigt at
være Videnskab og saa en Fare i, at Religionen blev henvist
til det rent Personliges og Subjektives Verden. Fra
modsat Side optraadte især Georg Brandes i forskellige
Bladartikler og senere i et lille Skrift „Dualismen i vor
nyeste Filosofi“. Jeg tog Del i Striden i et meget ubehjælpsomt
formet Skrift „Filosofi og Teologi“ (1866), der i
Hovedsagen var et Referat af den Diskussion, der efter
Fremkomsten af Martensens Dogmatik (1849) var bleven
ført om Forholdet mellem Filosofi og Teologi. Mit Forhold
til hele Sagen udtrykte jeg nogle Aar efter i mit
% PRINTER'S ERROR: "dennne" so printed (sense: denne); read at 600 dpi.
Doktorvita saaledes: „Jeg har altid opfattet dennne Anskuelse
[Rasmus Nielsens Opfattelse af Tro og Viden]
mere fra den skeptiske og negative Side, mere som et
Middel til en personlig Løsning af Vanskelighederne, end
som en egen positiv Lære.“ Det vil sige, jeg var mere
Kierkegaardianer end Nielsenianer. Brøchner gjorde senere
udtrykkeligt opmærksom paa Modsætningen mellem
Kierkegaards og Nielsens Standpunkt. Og denne Modsætning
% PRINTER'S ERROR: "( som" so printed, a full word-space after the
% opening parenthesis (sense: "(som"); read at 600 dpi.
blev ( som jeg har paavist i „Danske Filosofer“)
mere og mere udpræget under Nielsens fortsatte Udvikling
af sine Tanker, indtil han tilsidst (i sit Skrift „Om
Aandsdannelse“ 1877) helt fornegtede Kierkegaard og i
Grunden gav dennes teologiske Modstandere Ret.

For mit Vedkommende forholdt det sig saaledes, at jeg
var paa Vejen bort fra Teologien, og at det nu gjaldt om,
hvilken Opfattelse af Videnskab og af personlig Livsopfattelse
jeg kunde naa til. Der begyndte nu en ny Læreperiode
for mig. Det gjaldt at faa et Grundlag for min videre Udvikling
baade i videnskabelig og i rent personlig Henseende.
Dette var jeg mig naturligvis ikke bevidst dengang,
% --- p. 61 ---
\opage{61}men naar jeg ser tilbage finder jeg en uvilkaarlig Stræben
efter et Fodfæste. Der gik omtrent otte Aar, inden et
saadant Fodfæste blev naaet. Jeg maatte lære de forskellige
Hovedproblemers Ejendommelighed at kende og derved
forstaa baade deres Sammenhæng og deres Selvstændighed.
I et kort Udkast til en Forelæsning i Efteraaret
1874 over „Indledning til Filosofien“ indeholdtes Grundtrækkene
for Alt, hvad jeg senere mere eller mindre udførligt
har fremstillet.

Men jeg foregriber. --- Mine Studier af græsk Filosofi
samlede sig paa det Tidspunkt, jeg her staar ved, mere og
mere om en bestemt Opgave, som jeg behandlede i et
Skrift „Den antike Opfattelse af den menneskelige Villie“,
jeg i Aaret 1870 forsvarede for den filosofiske Doktorgrad.
Nogle Aar forud havde jeg besvaret Universitetets filosofiske
Prisopgave (for 1867---68) og vundet Guldmedaillen.
Ogsaa denne Opgave angik den menneskelige Villie,
idet den krævede en Undersøgelse af den Strid om Villiens
Frihed, der i sin Tid (1824 o. f.) blev ført i vor
Litteratur mellem Retslægen Howitz, Sibbern, A. S. Ørsted
o. fl. Jeg er ikke tilfreds med den Løsning, jeg antydede,
og Afhandlingen havde efter Censorernes (Nielsens
% PRINTER'S DEFECT: between "og" and "Brøchners" a raised space has
% inked (a small colon-like mark and a hairline, read at 600 dpi; the
% text layer reads "og;").  Not a character; not transcribed.
og Brøchners) Dom sin væsentlige Betydning ved den indgaaende
Karakteristik af de forskellige Standpunkter, der
havde gjort sig gældende, og af hvilke Howitz’ og Sibberns
især havde interesseret mig.

Jeg kom fra denne Tid af i nærmere Forhold til Brøchner.
Han raadede mig, da jeg havde vundet Guldmedaillen,
til at rejse udenlands, og min Fader forærede mig da en
Rejse. Jeg vaklede noget med Hensyn til, hvor jeg skulde
rejse hen. Jeg havde ved Siden af mit andet Arbejde søgt
at orientere mig med Hensyn til Filosofiens Udvikling efter
Hegels Tid i Tyskland, og her havde Gøttingerfilosofen
Hermann Lotze især tildraget sig min Opmærksomhed. Jeg
% SETTLED 23 Sep 2026: "tæn" is SUPPLIED by the editor and printed \supplied{tæn}: conjecture from the lost
% width and the sense (2nd and 3rd readers: probable; only the top of the k, a trace of the n, and "kte"
% survive).  "Men" in the next line: 2nd reader probable, 3rd reader certain.
% Transcriber's note: "tænkte" vs any other word in "-kte" -- SCAN DEFECT: the left
% end of the last two lines of p. 61 is blank white in the scan (paper
% loss or masking; nothing recoverable at 600 dpi).  Surviving: the top
% of the ascender of k, the top of a shoulder just before it (as of n),
% and "kte".  Measured at 600 dpi against "tænkte" on p. 55: there t-left
% to k-stem is 134 px; here the k stem stands at x=554, which puts the
% first letter at x~420, exactly the left margin (418--423).  So "tæn"
% fills the lost space; the letters themselves are not visible.
\supplied{tæn}kte mig Muligheden af at studere hos ham i Gøttingen.
% SCAN DEFECT (same patch): of the first letter only the lower half of
% a pointed V-vertex and a right stem survive -- the middle of a capital
% M, not the arches of m; the lost width (to x~420) fits M.  Read "Men".
Men jeg skønnede, at jeg kunde lære denne fine Tænker og
% p. 62 opens mid-sentence, not mid-word: p. 61 ends "denne fine Tænker og"
% --- p. 62 ---
\opage{62}aandrige Forfatter nok saa godt at kende gennem hans Bøger som gennem hans Forelæsninger, og et Ophold i en lille tysk Universitetsby vilde ikke kunne berige mig paa andre Omraader end det rent videnskabelige. Jeg mente, at jeg vilde have et mere alsidigt Udbytte af at tilbringe en Vinter i Paris, hvor der ikke blot kunde studeres, men hvor Kunst og Literatur og et pulserende Hovedstadsliv afgav Muligheder for mere alsidig Uddannelse. Jeg tilbragte en meget indholdsrig Vinter i Paris (1868---69), hørte Forelæsninger (bl. a. af Taine), skrev paa min Disputats, var en stadig Gæst paa Musæer og i Teatre og brugte mine Øjne til at iagttage Livet i dette Verdenscentrum. Her lærte jeg andre Tankeretninger at kende end dem, der stammede fra Tyskland, og som hidtil havde havt størst Indflydelse paa dem, der herhjemme gav sig af med Filosofi. Jeg blev især opmærksom paa den positivistiske Retning, repræsenteret dels ved Comte’s store Skikkelse, ved Siden af, og som Modsætning til, Hegel, den største systematiske Tænker i det nittende Aarhundrede, dels ved Herbert Spencers dristige, men storartede Generalisation af Udviklingsbegrebet. Det varede dog nogen Tid, før denne Retning fik en mere afgørende Betydning for mig. Jeg har stedse været langsom i Vendingen. Tankerne maatte stille modnes hos mig, før jeg kunde arbejde med dem.

% white line in the print, p. 62
\whiteline

Jeg har hidtil i denne Selvbiografi væsentligt støttet mig til min Erindring. Men for at give et mere umiddelbart og kildemæssigt Indtryk af den Periode i mit Liv, der i dette Afsnit er Tale om, vil jeg indskyde nogle Brudstykker af Breve til min Forlovede i Aarene 1866 til 1870, Breve, jeg fandt i hendes Gemmer efter hendes Død, og som --- i Forbindelse med de Breve fra hende, jeg havde opbevaret --- nu i mit 74ende Aar havde gjort mig det muligt at sætte mig tilbage i mine indre og ydre Kaar i hine Dage. Jeg forbigaar Familieforhold og Andet, som ikke er af Inter%
% --- p. 63 ---
\opage{63}esse for den Opgave, jeg har sat mig i dette Skrift. Brevstederne anføres i kronologisk Orden.

% white line in the print, p. 63, before the letters
\whiteline

30. Septbr. 1866. --- Prøver man paa at holde Livet gaaende paa idyllisk, beskuende Harmoni, saa vil man nok støde an paa Virkeligheden og erfare, at Forsoningen kun var i Tanken eller i Fantasien ... Det er derfor, jeg er bleven en Discipel af Rasmus Nielsen, fordi han saa stærkt fremhæver Modsætningen mellem Tro og Viden. Det er jo to helt forskellige Sjælstilstande. Troen er et „uforandret Sind“, men Viden hører til Beskuelsens Rige.

25. Novbr. 1866. --- Hvor mildt og roligt udbreder Livet sig ikke for os, naar vi tænke paa dem, i hvis Lod de store Kampe og Anfægtelser faldt ... „Alle de, som leve gudeligt i Jesu Christo, skulle forfølges.“ (Andet Brev til Timotheus) ... For at vort Livsvilkaar skal være vort Livs Sandhed, maa vi jo forstaa os selv overfor de store Forbilleder og ikke i Uklarhed lade, som om den hellige Grav var vel forvaret. Hvad der fordres, er Ydmyghed og Oprigtighed ... Saa vil Oprejsningen og den inderlige Selvbekræftelse ikke udeblive.

% date line without the em rule, so printed
1866 (Datoen mangler). Naar jeg sjældent gaar i Kirke, er det tildels, fordi jeg
% PRINTER'S ERROR: "ngen" so printed (sense: ingen); read at 600 dpi, no trace of an i
ngen Præst véd, der er oprigtig overfor den egentlige Kristendom. Der er Steder i det nye Testamente, man aldrig hører nogen Præst tale om. Det er ingenlunde, fordi jeg tror at kunne være mig selv nok, at jeg aldrig gaar i Kirke.

9. Decbr. 1866. --- Du sagde nyligt, at jeg sjældent omtalte dem derhjemme. Det kommer ikke af, at jeg ikke har levet med dem, jeg mener ikke blot i ydre Berørelse, men virkelig, indre Samleven. Tværtimod har jeg stedse følt en inderlig Glæde ved dem, og jeg tror, jeg forstaar dem langt bedre end tidligere. For et Par Aar siden var det egentligt lidt trangt for mig i den Henseende, som du nok véd og vel selv saa. Jeg følte mig ofte saa isoleret og forladt, saa forskellig fra de Andre. Men nu er jeg vis paa, at
% --- p. 64 ---
\opage{64}Skylden maa have ligget hos mig. At der er en Forskellighed, er sikkert nok; men hvad gør ogsaa det egentlig, naar bare det rette Øje er der til at se den? --- Jeg kunde være bleven en slem En, hvis jeg havde faaet Lov at gaa saadan som for et Par Aar siden. Men der var En, som hindrede det, som gjorde mit Væsen mere aabent og mildt og forjog den Kulde, som i lange Tider kunde hvile saa tungt over mig. Det er Dig, jeg skylder, at jeg igen med et frit og glad Sind kan gaa herhjemme mellem Forældre og Søskende, for hvem det dog vist ofte har været en Plage at se paa saadan en Dødbider, som jeg dengang var. Jeg mener slet ikke, at jeg er ganske udskreven fra Hospitalet, --- det bliver jeg vel aldrig, og det er vel i sin Orden; men vist er det, at jeg bliver bange for mig selv, naar jeg tænker paa, hvorledes jeg var for tre Aar siden ...

Jeg har i dette Efteraar lært Marie (min Søster) at kende. Sagen er vel den, at jeg har trængt mere til hende, da jeg ikke havde Dig. Hvor det dog er en ægte, uforfalsket, barnlig Sjæl! Der er Guld i hende, og det vil staa sin Prøve. Og Prøven vil maaske blive haard; hun kender jo kun til Medgang --- men ogsaa den er et Slags Prøve.

1. Marts 1867. --- Du sagde engang i Julen, at jeg aldrig skrev Noget om mine Studeringer ... Saavidt jeg kan komme til Klarhed med mig selv, er jeg overbevist om, at det er det ene Rigtige, jeg gør, naar jeg stræber efter en mere omfattende og dybere gaaende Erkendelse og Forstaaelse i filosofisk Henseende. Det \textit{maa} lede til et Maal, som for mig personlig vil være af stor Vigtighed; jeg føler saa sikkert, at mit Liv vilde blive holdningsløst og vaklende, naar jeg afviste disse Spørgsmaal og lod dem staa hen i en ubestemt Taage. I den Henseende nærer jeg ingen Tvivl, uagtet Maalet naturligvis ofte synes at skjule sig for mig, og jeg forekommer mig selv saa afmægtig og saa kummerlig. Et andet Spørgsmaal er det derimod, om det vil lykkes mig at komme til et saadant Resultat, at det ogsaa udadtil fik nogen Betydning. Man kan jo ikke komme længere, end Evnen strækker til. Men er Du ikke enig med mig i, at det
% --- p. 65 ---
\opage{65}første Hensyn er det vigtigste? (Derefter omtaltes de praktiske Udsigter; jeg tænker mig at tage Doktorgraden, virke som Lærer ved en Skole, og engang overtage en af de private Latinskoler her i Byen.) At blive Universitetslærer er noget, der i vort lille Land jo er et ligefremt Held, ganske uafset de Fordringer, det stiller til Begavelse og Personlighed.

7. April 1867. --- Nu skal jeg miste Skovgaard. Det kan du tro vil blive et stort Savn for mig. Han skal være Huslærer hos en Grev Juel-Vind-Frijs paa Lolland ... Saa skal jeg ikke høre hans kraftige Skridt paa Trappen og stærke Dundren paa Døren, som nu i saa mange Aar har bebudet mig et kært Besøg. Det er ingen liden Gavn at have En, man ret kan kalde sin Ven i Ordets fulde Betydning, med hvem man kan tale ud af Hjertet og føle, at man bliver ganske forstaaet.

16. Juni 1867. --- Du kalder mig „sand“. Jeg tror, at denne Sandhed hos mig kan blive til en Banghed, der maaske slapper mig ikke saa lidt ... Sand vilde jeg rigtignok fremfor Alt være. Denne Sandhed, som er Et med Inderlighed, er jo dog den eneste Betingelse for, at vi virkeligt kunne kalde noget vort ... Især er Usandheden saa vanskelig at komme efter, naar den skjuler sig bag et begejstrende Forsæt.

8. Septbr. 1867. --- Undertiden har jeg en Fornemmelse af, at jeg er paa en rigtig Vej --- det vil sige paa en Vej, som kan føre mig til, hvad jeg tror, jeg nærmest skulde opnaa.

15. Septbr. 1867. --- Det daglige Liv har sine Højdepunkter og sine store Fald i Henseende til Idealitet. Kunde man blot holde den Overbevisning fast, at det er efter Højdepunkterne, man skal dømme, og at det Skær, som fra dem falder hen over Livet, er det sande Lys, hvori dette skal ses.

22. Septbr. 1867. --- (Om Julius Petersen). Hvad jeg ofte har beundret som Energi og Villieskraft hos ham, tror jeg egentligt mere er en rastløs Higen, som formaar meget,
% --- p. 66 ---
\opage{66}naar Maalet synes at ligge nær ved Haanden, men som synker sammen i sig selv, naar det gælder taalmodig Udholdenhed. (Formedelst Skuffelse over at være forbigaaet ved Besættelse af en Reservelægepost, tænkte han paa at opgive at fortsætte sine Studier og at nedsætte sig i Provinsen.) Det vilde være smerteligt for mig at se En, til hvem jeg har havt saa mangen Fortrøstning, som jeg har havt til ham, give efter for Livets Tryk og Hindringer.

13. Oktbr. 1867. --- Hvad jeg stræber efter, hvad der mere og mere bliver mig en aandelig Velfærdssag, det er at arbejde mig frem til en fast Verdensanskuelse. Derpaa skal der naturligvis arbejdes gennem hele Livet, men jeg vilde dog gerne i min Ungdom have gennemført et saadant videnskabeligt Arbejde, at jeg deri kunde have et godt Grundlag ... Det er det, der staar for mig som \textit{mit} Livs Gerning.

27. Oktbr. 1867. --- Jeg tror altid, at det er en Idyl, der skal opføres i Livet, og derfor gaar det mig ofte saa ilde.

% p. 66: the comma after "Kulturformer" is ITALIC (3rd reading, measured: wider and more
% slanted than all 29 roman commas on the page); transcriber had it roman, 2nd reader undecidable.
7. Novbr. 1867. --- Det kan maaske interessere dig at høre, i hvilken Retning min Læsning hovedsagelig gaar. Da Spørgsmaalet om Tro og Viden, som pinte mig, da jeg studerede Teologi, og som siden stedse fremstiller sig for mig paany, førte mig til nærmere Beskæftigelse med Filosofiens Historie, blev jeg efterhaanden overtydet om, at uden at have opfattet den græske Tænkning i dens ejendommelige Hovedformer, vilde man ikke have den rette Forestilling om, hvad det egentligt drejer sig om. Grækerne havde nemlig naaet den højeste Fuldendelse af aandelig Udvikling, som kunde naas under rent menneskelige Vilkaar. Med Kristendommen kom Troen ind i Verden, og nye Kræfter sattes i Bevægelse, men Alt, hvad vi have af menneskelige \textit{Kulturformer,} det viser i sine Hovedtræk tilbage til Grækerne. Jeg har derfor allerede i et halvt Aars Tid beskæftiget mig med Platon og Aristoteles, og hvis det lykkes mig at gennemføre min Plan, vil jeg engang benytte et Emne herfra til en Doktordisputats. Ved Siden heraf
% --- p. 67 ---
\opage{67}sysler jeg med Historie og Islandsk, som jeg faar Brug for ved min Undervisning.

14. Novbr. 1867. --- I den sidste Tid har jeg læst endel højt af Shakespeare, især Søndag Aften, og ogsaa --- hvad der er en stor Nydelse --- hørt Poul Kierkegaard læse. For en 14 Dages Tid siden sade han, Troels Lund og jeg her en Aften og læste Shakespeare lige til Klokken blev næsten Et, uden at vi vidste af det.

24. Novbr. 1867. --- En Forandring, der er foregaaet med mig i dette Efteraar, bestaar i, at jeg ikke længere kan stemme med Rasmus Nielsen, til hvem jeg, som du véd, før idethele havde sluttet mig. Dels en Vaklen hos ham selv, dels en slaaende Kritik af Professor Brøchner, dels og fremfor Alt en fornyet Prøvelse af Sagen selv (ikke mindst i Samtaler med Poul Kierkegaard) har vist mig det Uholdbare i hans Anskuelse. Dog er jeg derved ikke ført nærmere til det Kirkelige ... Tværtimod er det med det samme gaaet op for mig, at jeg paa flere Punkter havde Meninger, som ikke havde nogen sand Rod hos mig, men hang fast som noget Udvortes og Overleveret. Kristendommen vil --- det er jeg fuldt og fast overbevist om --- altid være mig den højeste Sandhed, den, i hvilken jeg ene kan finde Livsgaaderne løste. Men kun hvad jeg personligt kan erfare og tilegne mig, kan være Sandhed for mig; og den Overbevisning er jeg kommen til, at hvad man mener sig at besidde, som gaar ud over det, beror paa ubevidst Hykleri. Den Anskuelse, hvis Uholdbarhed jeg nu er kommen til at indse, har naturligvis ikke været uden Betydning; det var et Udviklingstrin, som maatte gennemgaas.

18. Januar 1868. --- Jeg har næsten i hele det nye Aar været saa sløv og hypokonder, at jeg ofte hverken har vidst ud eller ind ... Det er aldrig noget udenfor mig, jeg trænger til i de Tilstande; det er Samling og Klaring, og navnlig Kraft til at holde mig selv under den tilbørlige Tugt, det skorter paa.

(Fra Aaret 1868 er der ikke mange Breve, da vi vare
% --- p. 68 ---
\opage{68}sammen i København fra Marts af, indtil jeg i Begyndelsen af December rejste til Udlandet.)

19. Juni 1869. --- (Om „den Periode, der ligger forud for min Forlovelsestid og endog tildels strakte sig ind i denne“). Jeg har undertiden kaldt denne Periode min sentimentale Periode, og med en vis Ret men rigtigere tror jeg dog, den betegnes ved at kaldes den asketiske Periode. Lykken og Glæden stode dengang for mig som uberettigede eller idetmindste mistænkelige Gæster ved Livets Bord. Bod og Forsagelse vare de Hovedtanker, hvorpaa Livsanskuelsen hvilede. Jeg siger ikke, at jeg dengang i Et og Alt levede som en Asket; men den tilsvarende Ængstelse var dog tilstede ... Havde jeg levet i et katolsk Land, da havde jeg, i hvert Tilfælde til enkelte Tider, ikke været langt fra Klosteret ... Er jeg bleven mere fri, uhildet og selvstændig, saa er jeg ogsaa bleven mere ydmyg end dengang.

11. Juli 1869. --- Den Dag idag har været mig god. Jeg har ret erfaret, hvilken Glæde det er at sysle med et Arbejde, som lægger Beslag paa Ens Interesse. Det er en aandrig Livsfornødenhed. Er det, man udretter, end nok saa beskedent, saa kommer der dog Øjeblikke, hvor man kan have den Bevidsthed: ogsaa jeg har følt de store Tanker, paa hvilke Tilværelsen hviler, og af hvilke den føres frem imod sit Maal. Er jeg end ikke blandt dem, hvem det er givet at gribe disse Tanker og stille dem frem i en ny og vækkende Form, saa har dog min Begejstring ikke været mindre. O, blot man engang ret kunde samle sig, ret koncentrere alle sine Kræfter og trænge enhver hæmmende Magt tilside! Maaske der dog var en Sten, som det kunde blive min Lod at føje til den store Bygning.

18. August 1869. --- Vi have begge, hver paa sin Maade, arbejdet os frem til større Klarhed og Fasthed ... Jeg har ofte med Forundring mærket, hvorledes Du saa Tingene paa samme Maade som jeg, og jeg tænker da især paa de friere religiøse Anskuelser, der vinde mere og mere Betydning og Sandhed for mig.

% --- p. 69 ---
% p. 69 opens with a new letter, indented
\opage{69}26. Juni 1870. --- Vi staa nu ved Enden af vor Forlovelsestid. Hvor meget er der dog ikke oplevet og gennemlevet i den Tid. Jeg tror egentligt, det er først nu, jeg er bleven ret selvstændig og altsaa en virkelig Mand ... Selvstændigheden bestaar i, at Tingene ikke længere oprulle sig for Tanken som et Panorama, men at man selv griber ind for at drage det Sande og Evige frem.

% white line in the print, p. 69, after the letters
\whiteline

Ifølge min Plan om at gaa ind i Skolevirksomhed søgte jeg gentagne Gange Adjunktpladser, og det var mig en Skuffelse, at det ikke lykkedes mig at faa en saadan Stilling, da jeg nu gerne vilde sætte Bo. Imidlertid vilde det vel da være gaaet saaledes, at jeg havde maattet opgive at fortsætte mine Studier i det Omfang, jeg havde tænkt mig. Og efter at jeg i Februar 1870 havde disputeret for Doktorgraden, søgte og fik jeg det Smith’ske Stipendium, som jeg kunde have Udsigt til at beholde nogle Aar, naar jeg fortsatte mine Studier. Paa det og paa Indtægten af min Skoleundervisning, samt et Tilskud fra min Fader, giftede jeg mig i Efteraaret 1870.

I Foraarssemestret 1871 begyndte jeg at læse som Privatdocent paa Universitetet. Min Bestræbelse for at orientere mig indenfor Samtidens Filosofi førte til Udgivelsen af to Bøger, „Filosofien i Tyskland efter Hegel“ (1872) og „Den engelske Filosofi i vore Dage“ (1874). Gennem kritisk Prøvelse af den engelske Filosofi, især som den repræsenteredes af Mill og Spencer, blev jeg ledet til at se Betydningen af en Tanke, der havde ligget bag ved Kierkegaards Filosoferen og endnu tidligere var bleven bevidst fremdragen af Kant, nemlig Tanken om det menneskelige Bevidsthedslivs samlende og enhedssøgende Karakter. Senere hentede jeg herfra den Hypotese, jeg lagde til Grund i min Psykologi. I Efteraaret 1874 holdt jeg en Forelæsning over „Indledning til Filosofien“, i hvilken jeg fra det saaledes vundne Synspunkt karakteriserede de forskellige filosofiske Problemer. Det blev fra denne Tid af efterhaan-%
% p. 69 ends mid-word: "efterhaan-" (printed hyphen kept); the word continues on p. 70
\gap{pp.~70--76}
% [text to be added: pp. 70--76]
% p. 77 opens mid-word: "overnat-/tede"; p. 76 not yet transcribed
% --- p. 77 ---
\opage{77}tede, for derefter næste Dag at vandre til mine Svigerforældres
Gaard, hvor jeg traf Hustru og Barn i bedste Velgaaende.

Enhver saadan Fodvandring, jeg har foretaget, har ladet mig leve mig
inderligere ind i vort Lands Natur, føle mig Et med det, baade paa dets
sandede Kyster, i de frodige Skove og mellem dets Bakker. Før og efter den
Tid, jeg her taler om, har jeg gennemvandret Nord-, Syd- og Vestsjælland,
Møen, Bornholm, det sydlige Fyn og store Strækninger i Jylland. Især den
jyske Natur har jeg, som jeg allerede har udtalt, Kærlighed til, og en
Sommerferie forekommer mig ikke normal, naar jeg ikke har fornyet
Bekendtskabet med den. Vendsyssel med dets vide Udsigter, dets Bakker, Enge
og Aaer, Midtjylland med Heder og Skove, Lyng og Moser, Vestjylland med
Klitter og Bølgebrus, Østjylland med Aarhusbugtens og Vejlefjords
Herligheder --- alt dette glæder mig ikke blot, mens jeg færdes der, men
Erindringsbilleder fra disse Egne kommer ofte igen i mine Fantasier. I Steen
Blichers Noveller er det ikke mindst hans prægtige Naturskildringer, der
have gjort ham til en af mine Yndlingsdigtere.

\whiteline

Naar jeg efter denne jyske Udflugt vender tilbage til mit Liv i København,
er der en Skikkelse, der rager frem for mig i Erindringen om Aarene omkring
1870. Det er Hans Brøchner. Jeg var kommen i Forbindelse med ham i Sommeren
1867, idet han, efter at have hørt mig tale som Opponent ved en Disputats,
sendte mig en venlig og opmuntrende Hilsen. Fra denne Tid af besøgte jeg ham
ofte. Jeg fik Raad af ham med Hensyn til mit Studium af Filosofiens
Historie, og da jeg havde skrevet en kritisk Anmeldelse af en af hans Bøger,
sendte han mig sine Modbemærkninger og bad mig komme en Aften og drøfte dem
med ham. I de følgende Aar, indtil hans Død (1875), besøgte jeg ham
regelmæssigt en Gang om Maaneden i hans ensomme Hjem, hvor han levede som
Enkemand med sine
% --- p. 78 ---
\opage{78}to Børn. Han omgikkes Ingen af sine Samtidige, med Undtagelse af
sin Ungdomsven Geheimeraad Vedel. Af Yngre kom foruden mig kun Georg Brandes
til ham, efterat hans nærmeste Discipel, den elskværdige og begavede
Kristian Møller var død. Naar man besøgte ham, var det som Regel ikke
filosofiske Emner, der kom paa Tale. Han følte i sin Ensomhed Lyst til at
høre Nyt fra den akademiske Verden og fra den øvrige Verden. Jeg ser ham
endnu sidde i sin Gyngestol eller gaa op og ned ad Gulvet med sine
karakteristiske lange og bestemte Skridt. Det usædvanligt smukke Ansigt
kunde da lyse af Velvillie og Interesse. Ofte fortalte han om sine Rejser
% PRINTER'S ERROR: "sine unge, tidligt afdøde Hustru" so printed (sense: sin)
eller om sine unge, tidligt afdøde Hustru. Om Costanza, den unge
Italienerinde, han havde været forelsket i i sin Ungdom, har han aldrig talt
til mig; det var først, da jeg udgav hans Breve, at jeg lærte denne Episode
af hans Liv at kende. I Brevene fremtræder den, saaledes som han skildrer
den, som en stemningsfuld Roman. Det synes at have været en Skriftefaders
Indgriben, der hindrede, at han kom til at føre Costanza op til Norden. Det
havde vist ogsaa været et voveligt Forsøg. --- Naar man rejste sig op og
vilde tage Afsked, blev han særligt livlig og havde mange Ting at spørge og
tale om; det var, som om han nødigt slap de faa Mellemled, der forbandt ham
med Omverdenen.

Han var en filosofisk Eneboer. Den Kreds af Tanker, han havde tilegnet sig,
var hans aandelige Hjem. Han havde ikke megen Evne til at meddele sig i fri
Form, skiftende sine Udtryk efter Situationer og Personligheder, og han
havde en Gang for alle sammenfattet sine Ideer i bestemte, meget abstrakte
Former. Men baade paa Forelæsninger og i Samtale kunde der komme en Vibreren
i Stemme og Ansigtsudtryk, et Tegn paa den dybe Stemning, der laa bagved og
vidnede om Tankens Sammenhæng med hele Personligheden. Han var en Tankens
Ridder, den sidste Repræsentant iblandt os for den Idealisme, der i
% PRINTER'S ERROR: "Tillidt" (Til-/lidt) so printed (sense: Tillid)
Tillidt til, at Tilværelsens Inderste kunde finde og havde fun%
% --- p. 79 ---
\opage{79}det Udtryk i Tankens Form, tillige i denne Form fandt
% the r of "Udtryk" (line 2 of p. 79) is broken in the type, but legible as r
Udtryk for det personlige Livs højeste Higen. Problemerne ere for dem, der
fulgte efter ham paa Filosofiens Omraade, blevne mere sammensatte og
mangfoldige. Paa en af sine sidste Forelæsninger erklærede Brøchner sig for
Tilhænger af „den idealistiske Filosofi, som Tyskland har Æren af at have
frembragt“, og der har neppe været mange, hos hvem denne Filosofis
% PRINTED AS IS: "Hovedtænker" (confirmed at the image, 500 dpi, 23 Sep 2026); the sense wants "Hovedtanker" (principal ideas), and the translation renders it so.
Hovedtænker i den Grad gennemtrængte og adlede Personligheden som hos ham.
Billedet af „den stille Professor“, som han er bleven kaldt, har fulgt mig
paa min Vej og mindet mig om, hvilken Værdi Tankelivet kan have, ogsaa ud
over det rent intellektuelle Omraade, naar der er Alvor i Tankens Arbejde.
Hans Venskab var mig en Ære og en Opmuntring, og det var med dyb Bevægelse,
jeg tog Afsked med ham paa hans Dødsleje. Det var Dagen, før han døde, jeg
saa ham sidste Gang. Han talte med fuldstændig Ro om sin Bortgang, og aftalte
med mig Ordningen af forskellige Sager. Da han bød mig Farvel og takkede mig
for mit Venskab, var jeg mere betagen end han. Hans Husbestyrerinde fortalte
senere, at han havde sagt ganske roligt, da jeg var gaaet: „Tag det Glas
bort, ti Dr. Høffdings Taarer ere deri.“

En lille Myte dannede der sig, da han var død, om hans sidste Øjeblikke. Han
døde den 17. December (1875) henad Eftermiddagen, altsaa i den mørkeste Tid.
Noget af det sidste, han sagde, var: „Vi maa have Lys herind.“ I den
nærmeste Tid efter hans Død hørte jeg fra flere Sider, at han skulde have
raabt paa Lys, og der var dem, der føjede til, at det nok ikke var saa
underligt, at den Mand raabte paa Lys! --- Myten synes dog at være død en
naturlig Død. Den fik en Tid nogen Næring ved, hvad der skete ved hans
Begravelse. Der var dengang ingen rent borgerlig Begravelse mulig. Han havde
ønsket at blive begravet fra Universitetet, og at en Kollega skulde tale;
men hvis det ikke kunde lade sig gøre, havde han nævnt en bestemt Præst, som
havde talt over hans unge Ven Kristian Møl%
% --- p. 80 ---
\opage{80}ler og ved den Lejlighed vist megen Takt. Men vedkommende Præst
synes at have misforstaaet Forholdet og lagde mere ind i den Afdødes meget
betingede Ønske, end berettiget var. Jeg tænkte paa at berigtige
Misforstaaelsen, men opgav det, da Brøchners ældste Ven, Geheimeraad Vedel,
bestemt fraraadede det.

I Fortalen til min Udgave af Brøchners og Molbechs Brevvexling og i „Danske
Filosofer“ har jeg givet en Karakteristik af Brøchner som Personlighed og som
Tænker, og mit Værk „Den nyere Filosofis Historie“ har jeg tilegnet „Mindet
om min Lærer og Ven“.

\secrule

Hvis jeg vilde prøve at skildre mit Hjemliv, efterat jeg havde sat Bo, vilde
jeg staa overfor den samme Vanskelighed, som jeg tidligere har omtalt, hvad
det angaar at danne sig et tydeligt Billede af sine Nærmeste. De faa Aar, i
hvilke dette Hjemliv bestod, staa for mig som en Enklave i Forhold til mit
øvrige Liv. Ikke som om det var en ren Idyl. Ved at se tilbage drager jeg
mig selv til Regnskab for mangen Time, hvor Mismod og Forstemthed, der ikke
altid havde dybtliggende Aarsager, gik ud over den, jeg allermindst vilde
bedrøve. Men der blev baaret over med mig, her, som i mit gamle Hjem. Og vi
udviklede os jevnsides, som det allerede var begyndt i Forlovelsesaarene. Min
Hustru var naturligvis ikke inde i det, som optog mig. Men naar det lykkedes
mig at finde et Udtryk for den menneskelige og personlige Side ved de Emner,
mine Studier angik, blev det grebet med Glæde og Forstaaelse. Jeg har altid
fundet det komisk, naar Hustruer uden videre havde deres Mænds Meninger,
især naar de, som ikke sjeldent ses, optræder mere dogmatisk og docerende end
deres mandlige Halvdele. Og jeg tør sige, at min Hustru havde en saadan Sans
for det Ægte og Selverhvervede, og tillige en saadan Takt, at hun kun gik
med til det, som hun ad sin egen Erfarings og Eftertankes Vej kunde tilegne
sig. Og i Poesiens Land kunde vi mødes som i en fælles Ver%
% --- p. 81 ---
\opage{81}den. Jeg læste for eller med hende Homer og Platon, Shakespeare og
Molière, foruden Samtidens Forfattere, i første Linie Bjørnson og Ibsen. Jeg
erindrer særligt, hvor optagen hun var af Homer. Mangen Aften, naar
Spørgsmaalet var, hvad vi skulde tage frem, sagde hun: „Lad os faa noget om
de gudlignende Helte!“ Vor Husandagt blev nu Poesi. I Forlovelsesaarene
havde det været Bibelen i
% p. 81: an ink speck at the foot of the "a" of "af" (layer reads ".af"); not a stop
Egenskab af dogmatisk Autoritet. Haand i Haand vare vi vandrede over fra
Dogmernes til Symbolernes og Poesiens Land. Ogsaa nu læste vi Bibelen, men
som stor Livspoesi.

Herlige Sommerrejser gjorde vi sammen, til Møen, til Aarhus og
Silkeborgegnen, eller vi mødtes paa hendes Forældres Gaard i Vendsyssel,
idet hun rejste forud, og jeg gerne gjorde en lille Udflugt, tildels tilfods,
og endte hos hende og hendes Familie. Jeg har omtalt Turen langs
Jammerbugten. En anden Gang tog jeg til Skagen. Det var i Sommeren 1875.
Skagen var dengang lige bleven opdaget af Drachmann. Ingen Hoteller og
Villaer førte endnu Krig med den frie Natur. Man kørte med Postvogn fra
Frederikshavn, tildels lige i Strandkanten, og fra Poststationen maatte man
vade i dybt Sand hen til den straatækte Kro. Her boede jeg næsten som eneste
Gæst; en Indremissionær boede der ogsaa. Sandets og Havets Verden kunde jeg
saaledes uforstyrret færdes i. Tilbageturen gjorde jeg tilvogns til
Frederikshavn, og derfra vandrede jeg den herlige Vej langs Stranden, med
Bangsbo Bakker paa højre Haand, ad Sæby til. Den lyse Sommereftermiddag
staar endnu klart i mit Minde, og ikke mindst det Øjeblik, da jeg lidt
udenfor Sæby mødte min Hustru, der var kørt mig imøde. ---

Vi omgikkes i København først og fremmest Julius Petersen og hans Hustru
Oliva, f. Munch. Han praktiserede nu paa Nørrebro og fik hurtigt en stor
Praxis, ikke mindst i Fattigkvartererne. Han var rastløst virksom, og
begyndte allerede nu sine Studier over Lægekunstens Historie. Naar han kom
hjem fra Sygebesøg i fattige Hjem, kunde han forarges over den Luxus, han
mente at se i de%
% --- p. 82 ---
\opage{82}res meget tarvelige Hjem, og han kunde sidde hensunken i
Forstemthed og Misfornøjelse. Men pludseligt kunde saa en helt anden Stemning
komme op, og ofte, naar vi troede, at Barometret stod lavt, kunde der komme
Indbydelse til at deltage i et lille Lag. Overfor hans Rastløshed og hans
mange Svingninger var der god Brug for Olivas Besindighed, klare Forstand og
faste Villie. Da Børneflokken hurtigt voxede, var der nok at gøre for hende.
Og stedse holdt hun sit Hus i god Gang, saa det endog blev et Samlingssted
for hans Venner. Her mødtes vi med Kristian Arentzen, Viggo Hørup og Holger
Drachmann og deres Hustruer, og i mange og lange Taler ved det festlige
Aftensbord drøftedes Tidens mange Spørgsmaal. --- Troels Lund havde jeg
gjort Bekendtskab med i Studenteraarene. Fra Filologien var han gaaet over
til Teologien, og efter sin Examen gav han sig, ligesom jeg, af med Studier
over græsk Filosofi og disputerede (1871) med en interessant Afhandling om
Sokrates. Senere gik han over til historiske Studier, foranlediget ved en
Ansættelse i Rigsarkivet (Geheimearchivet, som det dengang hed). Troels
Lunds fine, harmoniske Personlighed var det velgørende at omgaas, og hans
livlige Hustru gjorde sammen med ham deres Hjem til et Sted, hvor vi gerne
kom. Og paa den Tid kunde man ogsaa have den Fornøjelse at blive overrasket
ved et Besøg; det var saa hyggeligt, naar Lund og hans Hustru kom og slog
sig ned for Aftenen, tagende til Takke med, hvad Huset formaaede. --- Ogsaa
Heegaards Hjem vise mange gode Erindringer tilbage til. Ved min
Doktordisputats var han Opponent, og fra den Tid af kom vi til at staa
hinanden nær. Efterat have brudt med Rasmus Nielsen forholdt han sig
% p. 82: an ink speck at the foot of the "a" of "af" (layer reads ".af"); not a stop
søgende, optagen af ivrigt Studium af Kant og de nyere engelske Filosofer;
tillige beskæftigede han sig meget med Matematik. Da hans økonomiske Forhold
vare trange, maatte han paatage sig opslidende Bierhverv. Bl. A. skrev han
anonymt en Række Romaner, i hvilke han var stærkt paavirket af Dickens. De
fandt i nogle Aar
\gap{pp.~83--89}
% [text to be added: pp. 83--89]
% p. 90 opens mid-sentence: "filosofisk set var meget overfladisk"; p. 89 not yet transcribed
% --- p. 90 ---
\opage{90}filosofisk set var meget overfladisk, --- og at den havde
mere Blik for det, der maatte falde for Kritiken, end
for det, der bevarer sin Betydning, og for de nye Muligheder
af positivt Værd. For mit Vedkommende var
Hovedsagen, at jeg havde nok at gøre med at skaffe mig
fast Grund under Fødderne paa det Omraade, hvor jeg
personligt skulde arbejde. I Midten af 70’erne var jeg naat
til, hvad man kunde kalde en kritisk Positivisme (og hvad
jeg etsteds fra et andet Synspunkt har kaldt en kritisk Monisme).
Den stillede mig stadigt nye Opgaver, som jeg
ikke paa Forhaand havde kunnet se, men som gave mig
stadigt Arbejde. Jeg havde, som jeg har sagt i „Tilbageblikket“,
en Traad at spinde videre paa, hvad enten Tidsstrømningen
gik i den ene eller den anden Retning. Litteraturens
Svingninger, som moderne Litteraturhistorikere tro
at kunne følge fra Tiaar til Tiaar, saaledes at hvert Tiaar
faar sin Karakteristik, har jeg ikke gjort med, og jeg har
heller ikke hos den unge Slægt, jeg har havt med at gøre,
ret kunnet spore disse Rytmer. Men jeg er vel heller ikke
fin nok i mine Sanser.

Georg Brandes og jeg havde været Modstandere i Striden
om Tro og Viden. Senere traf vi ofte sammen i Troels
Lund’s Faders Hus. Da han sammen med sin Broder havde
grundlagt Tidsskriftet „det nittende Aarhundrede“, blev
jeg Medarbejder og kom derved i en nærmere og for mig
meget gavnlig Forbindelse med ham.

Han indbød mig en Aften, da han var bleven gift, og
han og hans Hustru besøgte ogsaa mit Hjem. Jeg har altid
fulgt ham med Beundring i hans litteraturkritiske Arbejder.
Han har for mig som for saa mange Andre aabnet
nye poetiske og litterære Verdener, og ved sin Evne til indgaaende
Karakteristik og Analyse har han kastet Lys over
saa mange af Litteraturens store Skikkelser. Jeg følte mig
i Naturel meget forskellig fra ham. Jeg tog tungere og
mere ubehjælpsomt paa Alting, og jeg manglede den Selvtillid,
der er Betingelsen for dristig Optræden. Hvad jeg
% --- p. 91 ---
\opage{91}har kunnet udrette, skyldes langt og ufortrødent Arbejde,
og mine Opgaver laa den store Offentlighed fjernere. Desuden
har jeg gennemgaaende havt den Opfattelse, at hvad
der er at ændre eller skaffe afvejen i Tankens eller i Samfundets
Verden, først kan erkendes, naar det Rette findes.
Det er vanskeligt at finde det Rette, men naar det findes,
bliver det klart, hvad det Falske er. Og de Former, i hvilke
Menneskers Trang til Støtte og Tilhold i Livet finder Hvile,
kunne kun ændres, naar selve denne Trang ændres; men
saadan Ændring foregaar meget langsomt, naar den skal
være grundig. Ogsaa her er det Sandheden, der viser, hvad
der er falsk, Positiviteten, der begrunder Negativiteten.
Som Spinoza siger: Sandheden vidner baade om sig selv
og om Usandheden (verum est index sui et falsi). Og
Comte har hævdet, at man kun fortrænger det, som man
erstatter. Denne Betragtningsmaade gjorde, at jeg stod
anderledes overfor religiøs og moralsk Polemik end Brandes.
Men hans Kritik er paa mange Maader kommen min
Virksomhed til Gode, idet den hos mange vakte Trang til
videregaaende Orientering paa Tankens Omraade. Jeg har
ganske vist i ikke faa Tilfælde set, hvorledes den Opfattelse
kom op, at det var nok at have raske og forsorne Meninger,
uden at man behøvede at bryde sig om, at Filosofien forlangte
en indtrængende Undersøgelse.

Hvad jeg fra den Tid, jeg her taler om, Halvfjerserne,
ikke mindst er Brandes taknemlig for, er de Raad, han gav
mig i Anledning af de Bidrag, jeg sendte til hans Tidsskrift.
Saaledes skrev han til mig den 7. Juni 1875: „Tillad
mig, ej saa meget som Redaktør som i Egenskab af sympatetisk
Medstræbende at rette et Par Spørgsmaal til Dem.
Har De nogensinde gjort det Experiment at læse, hvad De
giver i Trykken, op for en eller anden Mand, der ikke har
gaaet i en lærd Skole? Har De læst det højt for en Dame,
der ikke var særligt filosofisk dannet? Har De sat et Mærke
ved de Steder, han eller hun ikke forstod, og forsøgt at omskrive
dem med andre Ord? Jeg tør sige, at De ikke har
% --- p. 92 ---
\opage{92}gjort dette Forsøg. Men hvorfor ikke?“ Og derpaa anker
særligt over et Sted i en Artikel, jeg havde indsendt til
Tidsskriftet, og som unegteligt var meget lidet populært
udtrykt. Jeg fandt Anken rigtig; desværre rammer den vel
ogsaa meget af, hvad jeg senere har skrevet, uagtet jeg
virkeligt satte mig for at skrive saa almentilgængeligt som
muligt. Men der er nu engang Forskel paa Filosofi og Litteraturhistorie,
og desuden har det været min Lod at udvikle
mig langsomt og sent komme til Klarhed. Jeg havde kunnet
vente med at skrive, til alt kunde udtrykkes klart (om end
ikke for dem, der slet ingen Forudsætninger have); det
havde maaske ikke været til Skade for Literaturen. Men
jeg har altid selv lært meget ved at forsøge paa at fremstille
mine Tanker og min hidtil erhvervede Kundskab. Og
jeg havde i Brøchner et Exempel paa, hvad det kan føre til
at vente for længe med at skrive; han beklagede ofte, at
han havde villet vente med sine Udarbejdelser, til han fuldkomment
havde gennemarbejdet alle sine Emner, og da den
Tid syntes ham at være kommen, var han henved de Halvtreds,
og han havde levet sig helt ind i de abstrakte Udtryk,
der vare overleverede fra den hegelske Filosofi, og
som i og for sig kunde have gjort god Tjeneste som korte
Formler under hans ensomme Tænken. --- Jeg har dog
senere hen ofte været overrasket over, i hvor store Kredse
mange af mine Bøger er naaede ud.

Da jeg i Aaret 1876 havde udgivet min lille Bog „Om
Grundlaget for den humane Etik“, som jeg i anden Sammenhæng
skal komme tilbage til, skrev Brandes til mig
(24. Oktbr. 1876): „De er, skønt paa et andet Grundlag, den
„forklarede“ Brøchner, og det vilde visseligt have været
Brøchner en stor Glæde, om han havde oplevet denne Bog
af Dem. .... Jeg gratulerer Dem ogsaa til Formen i Deres
Bog. De har her vundet overordentlig i Frihed og Bevægelighed.
De sagde engang et mismodigt Ord til mig om Tørhed
i Deres Sprogform, hvor De var altfor streng imod
Dem selv .... Der er efter min Mening endnu for lidt
% --- p. 93 ---
\opage{93}Temperament i Deres Stil; De skriver endnu ikke tilstrækkeligt
med Deres hele Personlighed, det hele Menneske.
Men De er paa Vej til at gøre det.“ --- Iøvrigt minder
Brandes i dette Brev om Indskriften paa Velfærdskomiteens
Papirer: „Liberté, Egalité, Fraternité, ou la mort!“
--- og han mener, at jeg neppe nogensinde har været saa
ung, at jeg har kunnet sige dette. Deri har han ganske Ret.
Dog skurrede det lidt i mine Øren, naar han tilføjede: „De
maa være født til en vis Grad gammel.“ Om sig selv siger
han saa: „Jeg er bange for, at jeg engang har været meget
ung.“ ---

Det Forhold, hvori jeg var kommen til Brandes, var mig
til megen Glæde og Opmuntring. Efter Brøchners Død
følte jeg mig ensom, og da Brandes med sin store kunstneriske
Sans forenede Forstaaelse af filosofiske Undersøgelser,
var hans Breve mig stedse kærkomne. Nogle Aar efter
skrev han til mig fra Berlin: „Jeg glæder mig meget ved
at staa paa saa venskabelig Fod med Dem; der er Intet, jeg
hellere vilde end leve i broderlig Forstaaelse med alle Jevnaldrende
i Norden, der forske ærligt.“ Hvad der her udtaltes
saa varmt, vandt fuldstændig Genklang fra min Side.
Og denne Stemning forsvandt end ikke hos mig, da den
endel Aar senere blev sat paa Prøve ved en alvorlig og
skarp Modsætning mellem vore Anskuelser paa et væsentligt
Punkt.
% End of chapter II: the rest of p. 93 is blank.

% --- p. 94 ---
% ---- Chapter III.  Indhold: "III. 1874—1883 .... 94".  pp. 94--120 (PDF 96--122).
% Head read at the image: the bare numeral "III", centred, no full
% stop, no title; the page prints no folio.
\chapter*{III}
\addcontentsline{toc}{chapter}{III. 1874--1883}
\markboth{III. 1874--1883}{III. 1874--1883}

\opage{94}I Efteraaret 1874 holdt jeg en Forelæsning over „Indledning
til Filosofien“, i hvilken jeg i al Korthed fremstillede
den Opfattelse af de filosofiske Problemer og af Forholdet
mellem dem, som nu havde fæstnet sig hos mig. Og
derefter begyndte min egentlige Forfattervirksomhed, idet
jeg nu satte mig til Opgave at behandle de enkelte Problemer
hvert for sig. Det korte Program, jeg havde givet
i den nævnte Forelæsning, maatte naturligvis undergaa
mange Forandringer ved Udarbejdelsen i det Enkelte, og
der var stort Arbejde at gøre paa de enkelte Omraader. I
tre Perioder kan jeg inddele mit Arbejde, i de følgende 40
% p. 94: "Fra 1875 til 1887 drejede sig om" so printed, without "det" (confirmed at the 2nd and 3rd
% readings, normal spacing, no gap).  Not logged as a printer's error: the sentence can be construed.
Aar. Fra 1875 til 1887 drejede sig om Psykologi og Etik
i første Linie, fra 1887---1895 om Filosofiens Historie, og
efter 1895 er det Religionsfilosofi og Erkendelsesteori, der
især lægger Beslag paa mig. Saaledes som Filosofiens Stilling
endnu er, er det af nogen Betydning, at disse forskellige
Omraader, der paa saa mange Maader gribe over i
hverandre, kunne behandles af den Samme. Det er ikke
sagt, at det altid vil være muligt, og Filosofiens Problemer
ville da stille sig i vanskeligere Former. Vort Arbejde maa
vi til enhver Tid tage op, som Spørgsmaalene stille sig for
os; vi maa blot ikke glemme, at Spørgsmaal saavel som
Svar kan antage forskellig Karakter til forskellige Tider.
Jeg kan dog ikke tænke mig Filosofi uden en Stræben efter
at sammenfatte det menneskelige Tankearbejde under visse
almene Synspunkter. Og denne Stræben forberedes indenfor
de enkelte Videnskaber, jo mere disse selv søge at blive
% --- p. 95 ---
\opage{95}sig deres ledende Tanker bevidst. Der er da stedse en Opgave
for Filosofi som komparativ Videnskab.

De første Opgaver, der stillede sig for mig, vare: at
skrive en Etik ud fra rent menneskelige Forudsætninger,
og at udarbejde en Psykologi, ved hvilken Studiet af det
menneskelige Sjæleliv, der siden Sibberns Tid ikke havde
været dyrket herhjemme, kunde genoptages. Jeg følte fra
denne Tid af mere og mere min Virksomhed som en Fortsættelse
af Sibberns, eller rettere, jeg ønskede, at den kunde
betragtes saaledes. Saa meget jeg end skyldte Rasmus Nielsen
og Brøchner, saa jeg dog, at man maatte gaa andre
Veje, Veje, som „gamle Sibbern“ paa mange Strækninger
havde befaret, skønt en romantisk-spekulativ Baggrund paa
mange Punkter havde bestemt det Lys, i hvilket han saa
Spørgsmaalene.

I Efteraaret 1876 udgav jeg et lille Skrift „Om Grundlaget
for den humane Etik“. Jeg søgte her at vise, at der
gives et rent menneskeligt Grundlag for Livsførelse, som
er givet med Muligheden af, at Mennesket kan bøje sig
for en Lov, der peger ud over de enkelte Øjeblikke og ud
over de enkelte Individer. Skønt jeg tror, at jeg naaede et
Stykke frem i Klarhed med det Skrift, og at særligt Analysen
af Begrebet Autoritet endnu kan have nogen Interesse,
var der endnu en Ubestemthed tilbage, som hang sammen
med, at jeg endnu ikke skelnede mellem det Motiv, der fører
os til at kalde Handlinger gode eller onde, og som jeg
senere har kaldt Vurderingsmotivet, og det Motiv, som i
de enkelte specielle Tilfælde faar os til at fyldestgøre Vurderingsmotivets
Krav, og som jeg senere har kaldt Handlingsmotivet.
Da jeg fik fuld Klarhed over denne Forskel,
der saa ofte overses, kunde jeg skride til Udarbejdelsen af
en udførlig Etik. Men det skete først adskillige Aar senere.

Foreløbigt gik jeg op i psykologiske Studier og behandlede
fra Efteraaret 1876 af hvert Aar et eller andet Afsnit
af Psykologien paa mine Forelæsninger som Privatdocent.
Efterhaanden samlede der sig en anseelig Tilhørerkreds om
% p. 95 ends mid-sentence; p. 96 not yet transcribed

% ---- Chapter III.  Indhold: "III. 1874—1883 .... 94".  pp. 94--120 (PDF 96--122).
% (The head of chapter III, top of p. 94, stands above, inside the
% batch for pp. 90--95, which crosses the chapter boundary.)
\gap{pp.~96--120}
% [text to be added: pp. 96--120]

% ---- Chapter IV.  Indhold: "IV. 1883—1894 .... 121" (so read at the
% image; the text layer's "I2J" is a nick in the last figure).
% pp. 121--174 (PDF 123--176).
\chapter*{IV}
\addcontentsline{toc}{chapter}{IV. 1883--1894}
\markboth{IV. 1883--1894}{IV. 1883--1894}
\gap{Chapter IV, pp.~121--174,}
% [text to be added: pp. 121--174]

% ---- Chapter V.  Indhold: "V. 1894—1904 .... 175".  pp. 175--223 (PDF 177--225).
\chapter*{V}
\addcontentsline{toc}{chapter}{V. 1894--1904}
\markboth{V. 1894--1904}{V. 1894--1904}
\gap{Chapter V, pp.~175--223,}
% [text to be added: pp. 175--223]

% ---- Chapter VI.  Indhold: "VI· 1904—1914 .... 224".  pp. 224--265 (PDF 226--267).
% PRINTER'S ERROR: the head on p. 224 prints "V" -- the I did not ink,
% leaving a faint stub beside the V.  Transcribed as printed.
\chapter*{V}
\addcontentsline{toc}{chapter}{VI. 1904--1914}
\markboth{VI. 1904--1914}{VI. 1904--1914}
\gap{Chapter VI, pp.~224--265,}
% [text to be added: pp. 224--265]

% ---- Chapter VII.  Indhold: "VII. 1914—1922 .... 266".  pp. 266--310 (PDF 268--312).
\chapter*{VII}
\addcontentsline{toc}{chapter}{VII. 1914--1922}
\markboth{VII. 1914--1922}{VII. 1914--1922}
\gap{Chapter VII, pp.~266--310,}
% [text to be added: pp. 266--310]

% ---- Chapter VIII.  Indhold: "VIII. 1922—1924 .... 311".  pp. 311--324 (PDF 313--326).
\chapter*{VIII}
\addcontentsline{toc}{chapter}{VIII. 1922--1924}
\markboth{VIII. 1922--1924}{VIII. 1922--1924}
\gap{Chapter VIII, pp.~311--324,}
% [text to be added: pp. 311--324]

\backmatter

% ---- INDHOLD, PDF 327, unnumbered: p. [325] by sequence.  With the
% note beneath it on when the sections were written.  The fragment
% supplies the head and its table-of-contents entry.
% --- p. 325 ---
% INDHOLD, PDF 327, unnumbered: p. [325] by sequence (no folio printed;
% the page follows p. 324 = PDF 326).  Head INDHOLD letterspaced
% capitals -> \emph.  The list and the note are set in the same type,
% smaller than the body (line pitch 92 px at 600 dpi against 107--108 px
% on p. 59), so both are given \small here.  Roman numerals right-aligned
% in a column; year ranges with an em rule and no spaces; dot leaders;
% page numbers in old-style figures, all read at 600 dpi.
\chapter*{\emph{INDHOLD}}
\addcontentsline{toc}{chapter}{Indhold}
\markboth{Indhold}{Indhold}
\opage{325}%
\begin{center}
\begin{minipage}{0.85\textwidth}
\small
\hfill{\footnotesize Side}\par
\noindent\makebox[3em][r]{I.}\enspace 1843---1861\dotfill 7\par
\noindent\makebox[3em][r]{II.}\enspace 1861---1874\dotfill 38\par
\noindent\makebox[3em][r]{III.}\enspace 1874---1883\dotfill 94\par
% Page for IV read at 600 dpi as old-style 1, 2, 1.  The last figure has
% lost the left half of its foot serif, and a small blot of ink stands
% just below the baseline to its left; that is the text layer's "J".
% The figure is 121, and chapter IV's head does stand on p. 121.
\noindent\makebox[3em][r]{IV.}\enspace 1883---1894\dotfill 121\par
\noindent\makebox[3em][r]{V.}\enspace 1894---1904\dotfill 175\par
% PRINTER'S ERROR: "VI·" so printed -- the point after VI stands at
% mid-height (a raised point), not on the baseline as after every other
% numeral (sense: VI.).
\noindent\makebox[3em][r]{VI·}\enspace 1904---1914\dotfill 224\par
\noindent\makebox[3em][r]{VII.}\enspace 1914---1922\dotfill 266\par
\noindent\makebox[3em][r]{VIII.}\enspace 1922---1924\dotfill 311\par
\end{minipage}
\end{center}

\begin{small}
% The note: five lines, first line indented, same size as the list.
% Line 1 ends "1915" and line 2 begins "—16." (the range broken at the
% em rule); joined here without space, as the ranges are set elsewhere.
Afsnittene I.---VI. er nedskrevne i Vinteren 1915---16. --- Afsnit VII.
er skrevet i Aaret 1922. --- Afsnit VIII. i Foraaret 1924. --- Inden
% "Paranteser" so printed (the usual 1928 form is "Parenteser"; the
% word occurs nowhere else in the book).  Transcribed as printed; it may
% be the author's spelling rather than the compositor's.  A small stray
% mark stands above the P, not part of any letter.
Trykningen har et Gennemsyn ført til at tilføje nogle Paranteser og
Noter.
\end{small}

\secrule

\end{document}
